\section{AIME II 2015}
        \begin{enumerate}
        	\item (\textit{AIME II 2015}) Let $N$ be the least positive integer that is both $22$ percent less than one integer and $16$ percent greater than another integer. Find the remainder when $N$ is divided by $1000$.


 
 \textbf{Answer: }131



\textbf{Solution 1}

If $N$ is $22$ percent less than one integer $k$, then $N=\frac{78}{100}k=\frac{39}{50}k$. In addition, $N$ is $16$ percent greater than another integer $m$, so $N=\frac{116}{100}m=\frac{29}{25}m$. Therefore, $k$ is divisible by $50$ and $m$ is divisible by $25$. Setting these two equal, we have $\frac{39}{50}k=\frac{29}{25}m$. Multiplying by $50$ on both sides, we get $39k=58m$.

 The smallest integers $k$ and $m$ that satisfy this are $k=1450$ and $m=975$, so $N=1131$. The answer is $\boxed{131}$.

 





\textbf{Solution 2}

Continuing from Solution 1, we have $N=\frac{39}{50}k$ and $N=\frac{29}{25}m$. It follows that $k=\frac{50}{39}N$ and $m=\frac{25}{29}N$. Both $m$ and $k$ have to be integers, so, in order for that to be true, $N$ has to cancel the denominators of both $\frac{50}{39}$ and $\frac{25}{29}$. In other words, $N$ is a multiple of both $29$ and $39$. That makes $N=\operatorname{lcm}(29,39)=29\cdot39=1131$. The answer is $\boxed{131}$.

 





\textbf{Solution 3}

We can express $N$ as $0.78a$ and $1.16b$, where $a$ and $b$ are some positive integers. $N=0.78a=1.16b\implies100N=78a=116b.$ Let us try to find the smallest possible value of $100N$, first ignoring the integral constraint. 

 Obviously, we are just trying to find the LCM of $78$ and $116.$ They have no common factors but $2$, so we multiply $78$ and $116$ and divide by $2$ to get $4524.$ This is obviously not divisible by $100$, so this is not possible, as it would imply that $N=\dfrac{4524}{100},$ which is not an integer. This can be simplified to $\dfrac{1131}{25}$. 

 We know that any possible value of $N$ will be an integral multiple of this value; the smallest such $N$ is achieved by multiplying this value by $25.$ We arrive at $1131$, which is $\boxed{131}\mod1000.$ 

 


 




	\item (\textit{AIME II 2015}) In a new school, $40$ percent of the students are freshmen, $30$ percent are sophomores, $20$ percent are juniors, and $10$ percent are seniors. All freshmen are required to take Latin, and $80$ percent of sophomores, $50$ percent of the juniors, and $20$ percent of the seniors elect to take Latin. The probability that a randomly chosen Latin student is a sophomore is $\frac{m}{n}$, where $m$ and $n$ are relatively prime positive integers. Find $m+n$.


 
 \textbf{Answer: }025



\textbf{Solution 1}

We see that $40\% \cdot 100\% + 30\% \cdot 80\% + 20\% \cdot 50\% + 10\% \cdot 20\% = 76\%$ of students are learning Latin. In addition, $30\% \cdot 80\% = 24\%$ of students are sophomores learning Latin. Thus, our desired probability is $\dfrac{24}{76}=\dfrac{6}{19}$ and our answer is $6+19=\boxed{025}$.

 





\textbf{Solution 2}

Assume that there are $100$ students in the school. There are $40$ freshmen taking Latin, $24$ sophomores taking Latin, $10$ juniors taking Latin, and $2$ seniors taking Latin. We get the probability to be the number of sophomores taking Latin over the total number of students taking Latin, or $\dfrac{24}{76}$. Simplifying, we get $\dfrac{6}{19}$. Adding, we get $\boxed{025}$.

 


 




	\item (\textit{AIME II 2015}) Let $m$ be the least positive integer divisible by $17$ whose digits sum to $17$. Find $m$.


 
 \textbf{Answer: }476



\textbf{Solution 1}

The three-digit integers divisible by $17$, and their digit sum: 
 \[\begin{array}{c|c} m & s(m)\\ \hline 102 & 3 \\ 119 & 11\\ 136 & 10\\ 153 & 9\\ 170 & 8\\ 187 & 16\\ 204 & 6\\ 221 & 5\\ 238 & 13\\ 255 & 12\\ 272 & 11\\ 289 & 19\\ 306 & 9\\ 323 & 8\\ 340 & 7\\ 357 & 15\\ 374 & 14\\ 391 & 13\\ 408 & 12\\ 425 & 11\\ 442 & 10\\ 459 & 18\\ 476 & 17 \end{array}\]

 Thus the answer is $\boxed{476}$.

 





\textbf{Solution 2 (Shortcut)}

We can do the same thing as solution 1, except note the following fact: $102$ is a multiple of $17$ and its digits sum to $3$. 

 Therefore, we can add it onto an existing multiple of $17$ that we know of to have $s(m) = 14$, shown in the right-hand column, provided that its units digit is less than $8$ and its hundreds digit is less than $9$. Unfortunately, $68$ does not fit the criteria, but $374$ does, meaning that, instead of continually adding multiples of $17$, we can stop here and simply add $102$ to reach our final answer of $\boxed{476}$.

 (Comment from another person: Actually, this doesn't work because you can't be sure there are no numbers between 374 and 476 that work. This solution just lucks out.)

 





\textbf{Solution 3}

The digit sum of a base $10$ integer $m$ is just $m\pmod{9}$. In this problem, we know $17\mid m$, or $m=17k$ for a positive integer $k$.

 Also, we know that $m\equiv 17\equiv -1\pmod{9}$, or $17k\equiv -k\equiv -1\pmod{9}$.

 Obviously $k=1$ is a solution. This means in general, $k=9x+1$ is a solution for non-negative integer $x$.

 Checking the first few possible solutions, we find that $m=\boxed{476}$ is the first solution that has $s(m)=17$, and we're done.

 





\textbf{Solution 4}

Since the sum of the digits in the base-10 representation of $m$ is $17$, we must have $m\equiv 17 \pmod{9}$ or $m\equiv -1\pmod{9}$.We also know that since $m$ is divisible by 17, $m\equiv 0 \pmod{17}$.

 To solve this system of linear congruences, we can use the Chinese Remainder Theorem. If we set $m\equiv (-1)(17)(8)\pmod {153}$, we find that $m\equiv 0\pmod{17}$ and $m\equiv -1\pmod{9}$, because $17\cdot 8\equiv 136 \equiv 1\pmod{9}$. The trick to getting here was to find the number $x$ such that $17x\equiv 1\pmod{9}$, so that when we take things $\pmod{9}$, the $17$ goes away. We can do this using the Extended Euclidean Algorithm or by guess and check to find that $x\equiv 8\pmod{9}$.

 Finally, since $m\equiv 17\pmod{153}$, we repeatedly add multiples of $153$ until we get a number in which its digits sum to 17, which first happens when $m=\boxed{476}$.

 





\textbf{Solution 5}

We proceed by casework on the number of digits. Clearly the answer must have at least two digits, seeing as the maximum digit sum for a one-digit number is 9. The answer must also have less than 4 digits, because this is the AIME.

 Case 1: The answer is a 2-digit number.Represent the number as $10a + b$, where $0 < a \leq 9$ and $0 \leq b \leq 9$. The conditions of the problem translated into algebra are:
 \[17|10a+b\]
 \[a+b=17\]By the Euclidean Algorithm, this is equivalent to:
 \[17|9a\]9 is not a factor of 17, so $17|a$. So $a$ must be a multiple of 17, but this is impossible because of the conditions we placed on $a$ and $b$.(Alternatively, note that the only possible options are 89 and 98, and neither works.)

 Case 2: The answer is a 3-digit number.Represent the number as $100a+10b+c$, where $0 < a \leq 9$ and $0 \leq b,c \leq 9$. Translating the conditions again:
 \[17|100a+10b+c\]
 \[a+b+c=17\]
 \[17|99a+9b\]
 \[17|9(11a+b)\]
 \[17|11a+b\]Testing multiples of 17 yields $(4, 7, 6)$ as the minimal solution for $(a, b, c)$ and thus the answer is $\boxed{476}$.

 - gting

 Note: We have $11a+b \equiv \mod(17)$. Doing some experimentation for the smallest value of $a$, since we want to minimize $abc_{10}$, we get $11a+b=51 \implies a=4, b=7, c=6$. Trying $17$ and $34$ don't work, because $0 \leq c \leq 9$. Also, one can use long division to verify that $17$ is a divisor of $476$, specifically $17 \cdot 28=476$. So the answer is $476$.~First

 





\textbf{Solution 6 (Simple Casework)}

Let $m$ be the integer with digits $\overline{a_n a_{n-1} \dots a_1}$ such that $\sum a_i = 17$ and $17 \mid m$.

 For 1-digit numbers, $a_1 \le 9$, so no solution exists. For 2-digit numbers, $a_2 + a_1 = 17$ only allows $89$ and $98$, neither of which is divisible by $17$. Thus, $m$ must be at least a 3-digit integer.

 For a 3-digit number $m = 100a_3 + 10a_2 + a_1$, the condition $m \equiv 0 \pmod{17}$ simplifies to:
 \[100a_3 + 10a_2 + a_1 \equiv 15a_3 + 10a_2 + a_1 \equiv 0 \pmod{17}\]

 We test $a_3 = 1, 2, 3, 4, \dots$ to find the smallest $m$. For each $a_3$, we check if any value $\overline{a_2a_1} \in \{17k + 2a_3\}$ satisfies the digit sum $a_2 + a_1 = 17 - a_3$.

 Reference Multiples of 17: $\{00, 17, 34, 51, 68, 85\}$ (Realistically, you should do the following part in your head.)

 
                \begin{itemize}
                    \item Case $a_3 = 1$: $\overline{a_2a_1} \equiv 2 \pmod{17}$. Resulting set: $\{02, 19, 36, 53, 70, 87\}$. Target sum $S = 17 - 1 = 16$. (None work)\item Case $a_3 = 2$: $\overline{a_2a_1} \equiv 4 \pmod{17}$. Resulting set: $\{04, 21, 38, 55, 72, 89\}$. Target sum $S = 17 - 2 = 15$. (None work)\item Case $a_3 = 3$: $\overline{a_2a_1} \equiv 6 \pmod{17}$. Resulting set: $\{06, 23, 40, 57, 74, 91\}$. Target sum $S = 17 - 3 = 14$. (None work)\item Case $a_3 = 4$: $\overline{a_2a_1} \equiv 8 \pmod{17}$. Resulting set: $\{08, 25, 42, 59, 76, 93\}$. Target sum $S = 17 - 4 = 13$.
                \end{itemize}


                For 76, the sum $7+6=13$ matches $S$. 

 The smallest integer is $\boxed{476}$.

 


 




	\item (\textit{AIME II 2015}) In an isosceles trapezoid, the parallel bases have lengths $\log 3$ and $\log 192$, and the altitude to these bases has length $\log 16$. The perimeter of the trapezoid can be written in the form $\log 2^p 3^q$, where $p$ and $q$ are positive integers. Find $p + q$.


 
 \textbf{Answer: }018



\textbf{Solution}

Call the trapezoid $ABCD$ with $AB$ as the smaller base and $CD$ as the longer. Let the point where an altitude intersects the larger base be $E$, where $E$ is closer to $D$. 

 Subtract the two bases and divide to find that $ED$ is $\log 8$. The altitude can be expressed as $\frac{4}{3} \log 8$. Therefore, the two legs are $\frac{5}{3} \log 8$, or $\log 32$.

 The perimeter is thus $\log 32 + \log 32 + \log 192 + \log 3$ which is $\log 2^{16} 3^2$. So $p + q = \boxed{018}$

 


 





\textbf{Solution 2 (gratuitous wishful thinking)}

Set the base of the log as 2. Then call the trapezoid $ABCD$ with $CD$ as the longer base. Then have the two feet of the altitudes be $E$ and $F$, with $E$ and $F$ in position from left to right respectively. Then, $CF$ and $ED$ are $\log 192 - \log 3 = \log 64$ (from the log subtraction identity. Then $CF=EF=3$ (isosceles trapezoid and $\log 64$ being 6. Then the 2 legs of the trapezoid is $\sqrt{3^2+4^2}=5=\log 32$.

 And we have the answer:

 $\log 192 + \log 32 + \log 32 + \log 3 = \log(192 \cdot 32 \cdot 32 \cdot 3) = \log(2^6 \cdot 3 \cdot 2^5 \cdot 2^5 \cdot 3) = \log(2^{16} \cdot 3^2) \Rightarrow 16+2 = \boxed{18}$

 -dragoon

 





\textbf{Solution 3}

Let $ABCD$ be the trapezoid, where $\overline{AB} || \overline{CD}$ and $AB = \log 3$ and $CD = \log 192$. Draw altitudes from $A$ and $B$ to $\overline{CD}$ with feet at $E$ and $F$, respectively. $AB = \log 3$, so $EF = \log 3$. Now, we attempt to find $DE + FC$, or what's left of $CD$ after we take out $EF$. We make use of the two logarithmic rules:

 
 \[\log(xy) = \log x + \log y\]

 
 \[\log(x^a) = a\log(x)\]

 
 \[CD = \log 192 = \log (3 \cdot 2^6) = \log 3 + \log(2^6) = \log 3 + 6\log 2\]

 Thus, since $CD = DE + EF + FC = \log 3 + 6\log 2$, $CD - EF = \log 3 + 6\log 2 - \log 3 = 6\log 2 = DE + FC$.

 Now, why was finding $DE + FC$ important? Absolutely no reason! Just kidding, lol 🤣 Now, we essentially "glue" triangles $\triangle DAE$ and $\triangle BFC$ together to get $\triangle XC'D'$, where $X$ is the point where $A$ and $B$ became one. Note we can do this because $\triangle DAE$ and $\triangle BFC$ are both right triangles with a common leg length (the altitude of trapezoid $ABCD$).

 Triangle $XC'D'$ has a base of $C'D'$, which is just equal to $DE + FC = 6\log 2$. It is equal to $DE + FC$ because when we brought triangles $\triangle DAE$ and $\triangle BFC$ together, the length of $CD$ was not changed except for taking out $EF$.

 $XC' = XD'$ since $AD = BC$ because the problem tells us we have an isosceles trapezoid. Drop and altitude from $X$ to $C'D'$ The altitude has length $\log 16 = 4\log 2$. The altitude also bisects $C'D'$ since $\triangle XC'D'$ is isosceles. Let the foot of the altitude be $M$. Then $MD' = 3\log 2$ (Remember that C'D' was $6\log 2$, and then it got bisected by the altitude). Thus, the hypotenuse, $XD'$ must be $5\log 2$ from the Pythagorean Theorem or by noticing that you have a 3-4-5 right triangle with a similarity ratio of $\log 2$. Since $XD' = XC' = BC = AD$, $BC = AD = 5\log 2 = \log 2^5$. 

 Now, we have $CD = \log (3 \cdot 2^6)$, $AB = \log 3$, and $BC = AD = \log 2^5$. Thus, their sum is 

 
 \[\log (3 \cdot 2^6) + \log 3 + \log 2^5 + \log 2^5 = \log (2^16 \cdot 3^2)\]

 Thus, $p + q = 16 + 2 = \boxed{18}$. ~Extremelysupercooldude

 





\textbf{Solution 4}

Let $a=\log2$ and $b=\log3$ so that the base lengths are $\log3=b$ and $\log192=\log(3\cdot64)=\log3+\log\left(2^6\right)=6a+b$ and the altitudes are $\log16=4a$. Then we have the following picture: 

 
                    \begin{figure}[H]
                        \centering
                        \includegraphics[width=0.5\textwidth]{Solution/AIME_2015_II_4_Solution_Figure_1.png}
                    \end{figure}
                    

 Note that we have the two right triangles to the side; one of each of their bases is an altitude, which we know the length to be $4a$. The length of the other base can be calculated as $\dfrac{(6a+b)-b}2=3a$ via simple isosceles trapezoid geometry; it is clear that each right triangle is actually a $3-4-5$ triangle, so we know their hypotenuses (and the remaining unknown sides of our original trapezoid) have length $5a$ (because of the $3a$ and $4a$ bases). Our answer is therefore $5a+5a+b+6a+b=16a+2b=16\log2+2\log3=\log\left(2^{16}3^2\right)$; $p+q=16+2=\boxed{018}$. QED. 

 


 




	\item (\textit{AIME II 2015}) Two unit squares are selected at random without replacement from an $n \times n$ grid of unit squares. Find the least positive integer $n$ such that the probability that the two selected unit squares are horizontally or vertically adjacent is less than $\frac{1}{2015}$.


 
 \textbf{Answer: }090



\textbf{Solution 1}

Call the given grid "Grid A". Consider Grid B, where the vertices of Grid B fall in the centers of the squares of Grid A; thus, Grid B has dimensions $(n-1) \times (n-1)$. There is a one-to-one correspondence between the edges of Grid B and the number of adjacent pairs of unit squares in Grid A. The number of edges in Grid B is $2n(n-1)$, and the number of ways to pick two squares out of Grid A is $\dbinom{n^2}{2}$. So, the probability that the two chosen squares are adjacent is $\frac{2n(n-1)}{\binom{n^2}{2}} = \frac{2n(n-1)}{\frac{n^2(n^2-1)}{2}} = \frac{4}{n(n+1)}$. We wish to find the smallest positive integer $n$ such that $\frac{4}{n(n+1)} < \frac{1}{2015}$, and by inspection the first such $n$ is $\boxed{090}$.

 


 





\textbf{Solution 2}

Consider a $3 \times 3$ grid, where there are $4$ corner squares, $4$ edge squares, and $1$ center square. A $4 \times 4$ grid has $4$ corner squares, $8$ edge squares, and $4$ center squares. By examining simple cases, we can conclude that for a grid that is $n \times n$, there are always $4$ corners squares, $4(n-2)$ edge squares, and $n^2-4n+4$ center squares. 

 Each corner square is adjacent to $2$ other squares, edge squares to $3$ other squares, and center squares to $4$ other squares. In the problem, the second square can be any square that is not the first, which means there are $n^2-1$ to choose from. With this information, we can conclude that the probability that second unit square is adjacent to the first is $\frac{2}{n^2-1}(\frac{4}{n^2}) +\frac{3}{n^2-1}(\frac{4(n-2)}{n^2}) +\frac{4}{n^2-1}(\frac{n^2-4n+4}{n^2})$.

 Simplifying, we get $\frac{4}{n(n+1)}$ which we can set to be less than $\frac{1}{2015}$. By inspection, we find that the first such $n$ is $\boxed{090}$.

 


 





\textbf{Solution 3}

There are 3 cases in this problem. Number one, the center squares. Two, the edges, and three, the corners. The probability of getting one center square and an adjacent square is $\frac{(n-2)(n-2)}{n^2}$ multiplied by $\frac{4}{n^2 -1}$. Add that to the probability of an edge and an adjacent square( $\frac{4n-8}{n^2}$ multiplied by  $\frac{3}{n^2-1}$) and the probability of a corner and an adjacent square( $\frac{4}{n^2}$ multiplied by  $\frac{2}{n^2-1}$) to get  $\frac{4n^2-4n}{n^4-n^2}$. Simplify to get  $\frac{4}{n^2+n}$. With some experimentation, we realize that the smallest value of n is $\boxed{090}$.

 


 





\textbf{Solution 4 (cheese)}

Notice how a chosen unit square on the grid has 4 vertically & horizontally adjacent squares around it (not counting corners or sides.) That's $\frac{4}{n^2}$. Using this, we rewrite $\frac{1}{2015}$ as $\frac{4}{8060}$. Notice that the denominator $8060$ is really close to $90^2$, and the problem is asking for the least positive integer less than $\frac{1}{2015}$. Therefore, the closest possible estimation is $\boxed{090}$. We can check this by adding in our corners and sides. Easy multiplication and simplification finds us with $\boxed{090}$ as the correct answer.

 





\textbf{Solution 5}

Lets start with some simpler cases. For example, lets use a 2x2. We have $2$ horizontally and $2$ vertically, for a total of $4$. Then, our possible outcomes is just $\binom{4}{2}$, so our probability is just $\dfrac{2}{3}$. Now, lets do a 3x3. We have that there are $6$ horizontally and $6$ vertically, for a total of $12$. Our possible outcomes is just $\binom{9}{2} = 36$. Now, our probability is just $\dfrac{1}{3}$. Now, lets try to generalize. For each colum, there is $n-1$ ways, then we multiply that by $n$ for our $n$ columns. Then we multiply all that by $2$ for the rows to get $2n(n-1)$. For our denominator, we have $n^2$ squares, and we need to choose $2$ of them, so our probability is: 
 \[\dfrac{2n(n-1)}{\binom{n^2}{2}}\]. Lets simplify this. We have: 
 \[\dfrac{2n(n-1)}{\dfrac{n^2(n^2-1)}{2}} \Longrightarrow \dfrac{4}{n(n+1)}\]. So, we have this to be less than $\dfrac{1}{2015}$. So we have: 
 \[\dfrac{4}{n(n+1)} < \dfrac{1}{2015} \Longrightarrow n(n+1) > 8060\], so the smallest $n$ that satisfies this condition is $n = \boxed{090}$

 -jb2015007

 


 




	\item (\textit{AIME II 2015}) Steve says to Jon, "I am thinking of a polynomial whose roots are all positive integers. The polynomial has the form $P(x) = 2x^3-2ax^2+(a^2-81)x-c$ for some positive integers $a$ and $c$. Can you tell me the values of $a$ and $c$?"


 After some calculations, Jon says, "There is more than one such polynomial."


 Steve says, "You're right.  Here is the value of $a$." He writes down a positive integer and asks, "Can you tell me the value of $c$?"


 Jon says, "There are still two possible values of $c$."


 Find the sum of the two possible values of $c$.


 
 \textbf{Answer: }440



\textbf{Solution 1 (Algebra)}

We call the three roots (some may be equal to one another) $x_1$, $x_2$, and $x_3$.  Using Vieta's formulas, we get $x_1+x_2+x_3 = a$, $x_1 \cdot x_2+x_1 \cdot x_3+x_2 \cdot x_3 = \frac{a^2-81}{2}$, and $x_1 \cdot x_2 \cdot x_3 = \frac{c}{2}$. 

 Squaring our first equation we get $x_1^2+x_2^2+x_3^2+2 \cdot x_1 \cdot x_2+2 \cdot x_1 \cdot x_3+2 \cdot x_2 \cdot x_3 = a^2$.

 We can then subtract twice our second equation to get $x_1^2+x_2^2+x_3^2 = a^2-2 \cdot \frac{a^2-81}{2}$.

 Simplifying the right side:

 

 \begin{align*} a^2-2 \cdot \frac{a^2-81}{2} &= a^2-a^2+81\\ &= 81.\\ \end{align*}


 So, we know $x_1^2+x_2^2+x_3^2 = 81$.

 We can then list out all the triples of positive integers whose squares sum to $81$:

 We get $(1, 4, 8)$, $(3, 6, 6)$, and $(4, 4, 7)$.  

 These triples give $a$ values of $13$, $15$, and $15$, respectively, and $c$ values of $64$, $216$, and $224$, respectively.  

 We know that Jon still found two possible values of $c$ when Steve told him the $a$ value, so the $a$ value must be $15$.  Thus, the two $c$ values are $216$ and $224$, which sum to $\boxed{\text{440}}$.

 





\textbf{Solution 2 (Algebra + Brute Force)}

First things first. Vietas gives us the following:

 
 \begin{align} x_1+x_2+x_3 = a\\ x_1 \cdot x_2+x_1 \cdot x_3+x_2 \cdot x_3 = \frac{a^2-81}{2}\\ x_1 \cdot x_2 \cdot x_3 = \frac{c}{2} \end{align}

 From $(2)$, $a$ must have odd parity, meaning $a^2-81$ must be a multiple of $4$, which implies that both sides of $(2)$ are even. Then, from $(1)$, we see that an odd number of $x_1$, $x_2$, and $x_3$ must be odd, because we have already deduced that $a$ is odd. In order for both sides of $(2)$ to be even, there must only be one odd number and two even numbers.

 Now, the theoretical maximum value of the left side of $(2)$ is $3 \cdot \biggl(\frac{a}{3}\biggr)^2=\frac{a^2}{3}$. That means that the maximum bound of $a$ is where
 \[\frac{a^2}{3} > \frac{a^2-81}{2},\]which simplifies to $\sqrt{243} > a$, meaning
 \[16 > a.\]So now we have that $9<a$ from $(2)$, $a<16$, and $a$ is odd from $(2)$. This means that $a$ could equal $11$, $13$, or $15$. At this point, we have simplified the problem to the point where we can casework+ brute force. As said above, we arrive at our solutions of  $(1, 4, 8)$, $(3, 6, 6)$, and $(4, 4, 7)$, of which the last two return equal $a$ values. Then, $2(3 \cdot 6 \cdot 6+4 \cdot 4 \cdot 7)=\boxed{440}$ AWD.

 





\textbf{Solution 3 (Basic calculus)}

Since each of the roots is positive, the local maximum of the function must occur at a positive value of $x$. Taking $\frac{d}{dx}$ of the polynomial yields $6x^2-4ax+a^2-81$, which is equal to $0$ at the local maximum. Since this is a quadratic in $a$, we can find an expression for $a$ in terms of $x$. The quadratic formula gives $a=\frac{4x\pm\sqrt{324-8x^2}}{2}$, which simplifies to $a=2x\pm\sqrt{81-2x^2}$. We know that $a$ is a positive integer, and testing small positive integer values of $x$ yields $a=15$ or $a=1$ when $x=4$, and $a=15$ or $a=9$ when $x=6$. Because the value of $a$ alone does not determine the polynomial, $a$, $a$ must equal $15$.

 Now our polynomial equals $2x^3-30x^2+144x-c$. Because one root is less than (or equal to) the $x$ value at the local maximum (picture the graph of a cubic equation), it suffices to synthetically divide by small integer values of $x$ to see if the resulting quadratic also has positive integer roots. Dividing by $x=3$ leaves a quotient of $2x^2-24x+72=2(x-6)^2$, and dividing by $x=4$ leaves a quotient of $2x^2-22x+56=2(x-4)(x-7)$. Thus, $c=2\cdot 3\cdot 6\cdot 6=216$, or $c=2\cdot 4\cdot 4\cdot 7=224$. Our answer is $216+224=\boxed{440}$ ~bad\_at\_mathcounts

 


 


 




	\item (\textit{AIME II 2015}) Triangle $ABC$ has side lengths $AB = 12$, $BC = 25$, and $CA = 17$. Rectangle $PQRS$ has vertex $P$ on $\overline{AB}$, vertex $Q$ on $\overline{AC}$, and vertices $R$ and $S$ on $\overline{BC}$. In terms of the side length $PQ = w$, the area of $PQRS$ can be expressed as the quadratic polynomial


 
 \[\text{Area}(PQRS) = \alpha w - \beta \cdot w^2.\]


 Then the coefficient $\beta = \frac{m}{n}$, where $m$ and $n$ are relatively prime positive integers. Find $m+n$.


 
 \textbf{Answer: }161



\textbf{Solution 1}

If $\omega = 25$,  the area of rectangle $PQRS$ is $0$, so

 
 \[\alpha\omega - \beta\omega^2 = 25\alpha - 625\beta = 0\]

 and $\alpha = 25\beta$.  If $\omega = \frac{25}{2}$, we can reflect $APQ$ over $PQ$, $PBS$ over $PS$, and $QCR$ over $QR$ to completely cover rectangle $PQRS$, so the area of $PQRS$ is half the area of the triangle.  Using Heron's formula, since $s = \frac{12 + 17 + 25}{2} = 27$, 

 
 \[[ABC] = \sqrt{27 \cdot 15 \cdot 10 \cdot 2} = 90\]

 so

 
 \[45 = \alpha\omega - \beta\omega^2 = \frac{625}{2} \beta - \beta\frac{625}{4} = \beta\frac{625}{4}\]

 and 

 
 \[\beta = \frac{180}{625} = \frac{36}{125}\]

 so the answer is $m + n = 36 + 125 = \boxed{161}$.

 





\textbf{Solution 2}


                    \begin{figure}[H]
                        \centering
                        \includegraphics[width=0.5\textwidth]{Solution/AIME_2015_II_7_Solution_Figure_1.png}
                    \end{figure}
                    

 Similar triangles can also solve the problem.

 First, solve for the area of the triangle. $[ABC] = 90$. This can be done by Heron's Formula or placing an $8-15-17$ right triangle on $AC$ and solving. (The $8$ side would be collinear with line $AB$)

 After finding the area, solve for the altitude to $BC$. Let $E$ be the intersection of the altitude from $A$ and side $BC$. Then $AE = \frac{36}{5}$. Solving for $BE$ using the Pythagorean Formula, we get $BE = \frac{48}{5}$. We then know that $CE = \frac{77}{5}$.

 Now consider the rectangle $PQRS$. Since $SR$ is collinear with $BC$ and parallel to $PQ$, $PQ$ is parallel to $BC$ meaning $\Delta APQ$ is similar to $\Delta ABC$. 

 Let $F$ be the intersection between $AE$ and $PQ$. By the similar triangles, we know that $\frac{PF}{FQ}=\frac{BE}{EC} = \frac{48}{77}$. Since $PF+FQ=PQ=\omega$. We can solve for $PF$ and $FQ$ in terms of $\omega$. We get that $PF=\frac{48}{125} \omega$ and $FQ=\frac{77}{125} \omega$.

 Let's work with $PF$. We know that $PQ$ is parallel to $BC$ so $\Delta APF$ is similar to $\Delta ABE$. We can set up the proportion:

 $\frac{AF}{PF}=\frac{AE}{BE}=\frac{3}{4}$. Solving for $AF$, $AF = \frac{3}{4} PF = \frac{3}{4} \cdot \frac{48}{125} \omega = \frac{36}{125} \omega$. 

 We can solve for $PS$ then since we know that $PS=FE$ and $FE= AE - AF = \frac{36}{5} - \frac{36}{125} \omega$.

 Therefore, $[PQRS] = PQ \cdot PS = \omega (\frac{36}{5} - \frac{36}{125} \omega) = \frac{36}{5}\omega - \frac{36}{125} \omega^2$.

 This means that $\beta = \frac{36}{125} \Rightarrow (m,n) = (36,125) \Rightarrow m+n = \boxed{161}$.

 





\textbf{Solution 3}

Heron's Formula gives $[ABC] = \sqrt{27 \cdot 15 \cdot 10 \cdot 2} = 90,$ so the altitude from $A$ to $BC$ has length $\dfrac{2[ABC]}{BC} = \dfrac{36}{5}.$

 Now, draw a parallel to $AB$ from $Q$, intersecting $BC$ at $T$. Then $BT = w$ in parallelogram $QPBT$, and so $CT = 25 - w$. Clearly, $CQT$ and $CAB$ are similar triangles, and so their altitudes have lengths proportional to their corresponding base sides, and so
 \[\frac{QR}{\frac{36}{5}} = \frac{25 - w}{25}.\]Solving gives $[PQRS] = \dfrac{36}{5} \cdot \dfrac{25 - w}{25} \cdot w = \dfrac{36w}{5} - \dfrac{36w^2}{125}$, so the answer is $36 + 125 = 161$.

 





\textbf{Solution 4}

Using the diagram from Solution 2 above, label $AF$ to be $h$. Through Heron's formula, the area of $\triangle ABC$ turns out to be $90$, so using $AE$ as the height and $BC$ as the base yields $AE=\frac{36}{5}$. Now, through the use of similarity between $\triangle APQ$ and $\triangle ABC$, you find $\frac{w}{25}=\frac{h}{36/5}$. Thus, $h=\frac{36w}{125}$. To find the height of the rectangle, subtract $h$ from $\frac{36}{5}$ to get $\left(\frac{36}{5}-\frac{36w}{125}\right)$, and multiply this by the other given side $w$ to get $\frac{36w}{5}-\frac{36w^2}{125}$ for the area of the rectangle. Finally, $36+125=\boxed{161}$.

 





\textbf{Solution 5}

Using the diagram as shown in Solution 2, let $AE=h$ and $AP=L$Now, by Heron's formula, we find that the $[ABC]=90$. Hence, $h=\frac{36}{5}$

 Now, we see that $\sin{B}=\frac{PS}{12-L}\implies PS=\sin{B}(12-L)$We easily find that $\sin{B}=\frac{3}{5}$ 

 Hence, $PS=\frac{3}{5}(12-L)$

 Now, we see that $[PQRS]=\frac{3}{5}(12-L)(w)$

 Now, it is obvious that we want to find $L$ in terms of $W$.

 Looking at the diagram, we see that because $PQRS$ is a rectangle, $\triangle{APQ}\sim{\triangle{ABC}}$

 Hence.. we can now set up similar triangles.

 We have that $\frac{AP}{AB}=\frac{PQ}{BC}\implies \frac{L}{12}=\frac{W}{25}\implies 25L=12W\implies L=\frac{12W}{25}$.

 Plugging back in..

 
$[PQRS]=\frac{3w}{5}(12-(\frac{12W}{25}))\implies \frac{3w}{5}(\frac{300-12W}{25})\implies \frac{900W-36W^2}{125}$

 Simplifying, we get $\frac{36W}{5}-\frac{36W^2}{125}$

 Hence, $125+36=\boxed{161}$

 





\textbf{Solution 6}

Proceed as in solution 1. When $\omega$ is equal to zero, $\alpha - \beta\omega=\alpha$ is equal to the altitude. This means that $25\beta$ is equal to $\frac{36}{5}$, so $\beta = \frac{36}{125}$, yielding $\boxed{161}$.

 





\textbf{Solution 7(Quick and Easy)}

Using the diagram in Solution 2, the area of $PQRS$ is ($w$)($PR$), we only need to find $PR$. Because $ABC$ and $APQ$ are similar in a ratio of $25:w.$ Because of $AB$:$AC$ $=12:15$ and $AP$:$AB$= $w:25,$ we can derive that $BP=\frac{(25-w)(12)}{25}.$ By using LoC on $ABC$, it is obvious that $cos(B)=\frac{4}{5}$ and $sin(B)=\frac{3}{5}.$ $PR=\frac{3bp}{5}=\frac{(25-w)(36)}{125}.$ Multiply $w$, we get $\frac{36W}{5}-\frac{36W^2}{125}$ which means the answer is  $\boxed{161}$

 fixed latex errors, ~eqb5000

 





\textbf{Solution 7 (fake solve)}

We are given that the area of rectangle $PQRS$ is a quadratic polynomial in terms of its width $\omega = PQ$. Let $h$ be the height of the rectangle (the length of $PS$ and $QR$). The area is given by $Area(\omega) = \omega h$.

 \textbf{Step 1: Identify the roots of the quadratic} A quadratic is uniquely determined by three points. We can identify two "degenerate" cases where the area of the rectangle is zero:

 
            \begin{enumerate}
                \item If $\omega = 0$, the rectangle collapses into a vertical line segment (the altitude), so $Area(0) = 0$.\item If $\omega = 25$ (the length of base $BC$), the rectangle collapses into the segment $BC$ itself, meaning its height $h = 0$, so $Area(25) = 0$.
            \end{enumerate}
            Since the roots are $0$ and $25$, the area polynomial must be of the form:

 
            \begin{enumerate}[label=(\roman*)]
                \item $Area(\omega) = k\omega(\omega - 25) = k\omega^2 - 25k\omega$
            \end{enumerate}
            Comparing this to the given form $Area(\omega) = \alpha \omega - \beta \omega^2$, we see that $\beta = -k$.

 \textbf{Step 2: Find the altitude of triangle ABC} To find the constant $k$, we need one more non-zero point. First, we calculate the area of $\triangle ABC$ using Heron's Formula. The semi-perimeter is $s = \frac{12+25+17}{2} = 27$.

 
            \begin{enumerate}[label=(\roman*)]
                \item $Area(\triangle ABC) = \sqrt{27(27-12)(27-25)(27-17)} = \sqrt{27 \cdot 15 \cdot 2 \cdot 10} = 90$
            \end{enumerate}
            Let $H$ be the altitude from $A$ to $BC$. Using $Area = \frac{1}{2}bh$:

 
            \begin{enumerate}[label=(\roman*)]
                \item $90 = \frac{1}{2}(25)(H) \implies H = \frac{180}{25} = \frac{36}{5}$
            \end{enumerate}
            \textbf{Step 3: Use similar triangles for a test point} Consider the case where $\omega = \frac{25}{2}$ (the width is half the base). By similar triangles ($\triangle APQ \sim \triangle ABC$), the height of the small triangle above the rectangle must be half of the total altitude $H$. Thus, the height of the rectangle $h$ is also half the total altitude:

 
            \begin{enumerate}[label=(\roman*)]
                \item $h = \frac{1}{2}H = \frac{18}{5}$
            \end{enumerate}
            Our test point is $(\omega, Area) = \left(\frac{25}{2}, \frac{25}{2} \cdot \frac{18}{5}\right) = \left(\frac{25}{2}, 45\right)$.

 \textbf{Step 4: Solve for beta} Substitute the point into our quadratic equation:

 
            \begin{enumerate}[label=(\roman*)]
                \item $45 = k\left(\frac{25}{2}\right)\left(\frac{25}{2} - 25\right)$\item $45 = k\left(\frac{25}{2}\right)\left(-\frac{25}{2}\right)$\item $45 = -\frac{625}{4}k$\item $k = -\frac{180}{625} = -\frac{36}{125}$
            \end{enumerate}
            Since $\beta = -k$, we have $\beta = \frac{36}{125}$. The values $m=36$ and $n=125$ are relatively prime, so $m+n = 36 + 125 = \boxed{161}$.

 


 




	\item (\textit{AIME II 2015}) Let $a$ and $b$ be positive integers satisfying $\frac{ab+1}{a+b} < \frac{3}{2}$. The maximum possible value of $\frac{a^3b^3+1}{a^3+b^3}$ is $\frac{p}{q}$, where $p$ and $q$ are relatively prime positive integers. Find $p+q$.


 
 \textbf{Answer: }036



\textbf{Solution 1}

Let us call the quantity $\frac{a^3b^3+1}{a^3+b^3}$ as $N$ for convenience. Knowing that $a$ and $b$ are positive integers, we can legitimately rearrange the given inequality so that $a$ is by itself, which makes it easier to determine the pairs of $(a, b)$ that work. Doing so, we have 
 \[\frac{ab+1}{a+b} < \frac{3}{2}\] 
 \[\implies 2ab + 2 < 3a + 3b \implies 2ab - 3a < 3b - 2\] 
 \[\implies a < \frac{3b - 2}{2b - 3}.\] Now, observe that if $b = 1$ we have that $N = \frac{a^3 + 1}{a^3 + 1} = 1$, regardless of the value of $a$. If $a = 1$, we have the same result: that $N = \frac{b^3 + 1}{b^3 + 1} = 1$, regardless of the value of $b$. Hence, we want to find pairs of positive integers $(a, b)$ existing such that neither $a$ nor $b$ is equal to $1$, and that the conditions given in the problem are satisfied in order to check that the maximum value for $N$ is not $1$.

 
To avoid the possibility that $a = 1$, we want to find values of $b$ such that $\frac{3b - 2}{2b - 3} > 2$. If we do this, we will have that $a < \frac{3b - 2}{2b - 3} = k$, where $k$ is greater than $2$, and this allows us to choose values of $a$ greater than $1$. Again, since $b$ is a positive integer, and we want $b > 1$, we can legitimately multiply both sides of $\frac{3b - 2}{2b - 3} > 2$ by $2b - 3$ to get $3b - 2 > 4b - 6 \implies b < 4$. For $b = 3$, we have that $a < \frac{7}{3}$, so the only possibility for $a$ greater than $1$ is obviously $2$. Plugging these values into $N$, we have that $N = \frac{8(27) + 1}{8 + 27} = \frac{217}{35} = \frac{31}{5}$. For $b = 2$, we have that $a < \frac{4}{1} = 4$. Plugging $a = 3$ and $b = 2$ in for $N$ yields the same result of $\frac{31}{5}$, but plugging $a = 2$ and $b = 2$ into $N$ yields that $N = \frac{8(8) + 1}{8 + 8} = \frac{65}{16}$. Clearly, $\frac{31}{5}$ is the largest value we can have for $N$, so our answer is $31 + 5 = \boxed{036}$.

 (Technically, we would have to find that b > 1 before dividing both sides of the inequality by 2b - 3, but otherwise this solution is correct)

 





\textbf{Solution 2 (nifty solution)}


 \[\frac{ab + 1}{a + b} < \frac{3}{2} \rightarrow 2ab + 2 < 3a + 3b \rightarrow 4ab - 6a - 6b + 4 < 0 \rightarrow (2a - 3)(2b - 3) < 5.\]

 
 \[2a - 3, 2b - 3 \in \{x \neq 2k, k \in Z \} \rightarrow\] $(2a-3)$ and $(2b-3)$ each cannot be even or else $a$ and $b$ will not be integers
 \[(2a - 3)(2b - 3) = 1, 3 \rightarrow (2a - 3, 2b - 3) = (1, 1), (1, 3), (3, 1).\]
 \[(a, b) = (2, 2), (2, 3), (3, 2).\]
 \[\frac{a^3 b^3 + 1}{a^3 + b^3} = \frac{65}{16}, \frac{31}{5}.\]
 \[\frac{31}{5} \rightarrow \boxed{036}.\]

 ($a=1 \rightarrow b=1, 2, 3... \rightarrow \frac{b^3+1}{b^3+1}=1$; $1<31/5$).

 





\textbf{Solution 3}

Notice that for $\frac{a^3b^3+1}{a^3+b^3}$ to be maximized, $\frac{ab+1}{a+b}$ has to be maximized. We simplify as above to $2ab + 2 < 3a + 3b$, which is $(a-\frac{3}{2})(b-\frac{3}{2}) < \frac{5}{4}$. To maximize, $a$ has to be as close to $b$ as possible, making $a$ close to $\frac{3+\sqrt{5}}{2}$. Because $a$ and $b$ are positive integers, $a = 3$, and checking back gives $b = 2$ as the maximum or the other way around, which the answer is thus $\frac{216+1}{27+8} = \frac{217}{35} = \frac{31}{5} \rightarrow \boxed{036}$.

 





\textbf{Solution 4}

Guess and check a few values of $a$ and $b$ you will find two things. One, that the highest values of $a$ and $b$ are the closest to $\frac{3}{2}$. Two, that the pair $(2,3)$ is the highest possible value of $a, b$. So plugging in $a=2$ and $b=3$ we get $((2*3)^3)+1/(8+27)$ = $\frac{217}{35}$ = $\frac{31}{5}$, and $31+5 = 36$.Solution by Helloitsaaryan

 





\textbf{Solution 5}

WLOG, assume $a \geq b$. We solve for $a$ from the original equation to get $a < \frac{3b - 2}{2b - 3}$. Since $b \leq a$, we can solve $b < \frac{3b - 2}{2b - 3}$. Solving, we get $\frac{3 - \sqrt{5}}{2} < b < \frac{3 + \sqrt{5}}{2}$. Approximating values, we get $1 \leq b < 3$. Thus, $b = 1, 2$. Now, plugging back $a < \frac{3b - 2}{2b - 3}$, we see that $b = 1$ doesn't work. Thus, we plug in $b = 2$ to get $2 \leq a < 4$. Thus, $a = 2, 3$. Clearly, we can plug in $(a, b) = (2, 2), (3, 2)$ and see that $(3, 2)$ gives you the maximum of $\frac{31}{5}$. Thus, the answer is $\boxed{036}$.

 





\textbf{Solution 6}

Firstly, if we assume that $a=1$, we see that $\frac{a^{3}b^{3}+1}{a^{3}+b^{3}} = \frac{b^{3}+1}{1+b^{3}} = 1$, which is not a fraction. Therefore, we know that neither $a$ nor $b$ can equal 1.

 Starting with $\frac{ab+1}{a+b}<\frac{3}{2}$, cross multiply and move the variables to the right to obtain $2ab-3a-3b<-2$. After multiplying both sides by 2, Simon's favorite factoring trick (SFFT) is applicable, giving us $(2a-3)(2b-3)<5$. Now, we assume $a=2$. Then we see that $b$ can equal 2 or 3. If $a=3$, the only possible $b$ value is 2, and there are no other possible values for $b$ if $a>3$. So the possible ordered pairs for $(a,b)$ are $(2,2), (2,3)$, and $(3,2)$. It is easily noticeable that $(2,3)$ and $(3,2)$ both produce $\frac{a^{3}b^{3}+1}{a^{3}+b^{3}}=\frac{31}{5}$, which is the maximum value. We then add $31+5=\fbox{036}$

 





\textbf{Solution 7 (Rigor Graph Analysis)}


                    \begin{figure}[H]
                        \centering
                        \includegraphics[width=0.5\textwidth]{Solution/AIME_2015_II_8_Solution_Figure_1.png}
                    \end{figure}
                    Define $x = a+b$ and $y = ab$. Then, the first inequality becomes $\frac{y+1}{x} < \frac{3}{2}$ so $3x > 2y+2$, highlighted blue, and through the AM-GM inequality, $x^2 \ge 4y$, highlighted red.

 For maximality to be reached, the points must lie on or near the quadratic, while in the blue region and red region overlap. Therefore, we can test points $(x,y) = (1,0), (2,1), (3,2), (4,4), (5,6)$. Note that every other point lies farther away from the quadratic starting $x = 6$, so no use in checking them. Also, the line $3x=2y+2$ only intersects the quadratic at two points, and so afterward, values of $x$ and $y$ force the inequality too minor. We can manually plug and chug the values $x$ and $y$ into the requested result $\frac{y^3 + 1}{x(x^2 - 3y)}$, and see that the maximum occurs at $(5,6)$ (duh...). So, we have $\frac{217}{35} = \frac{31}{5}$, which occurs when $a = 2,3$ and $b=3,2$ respectively. The answer is $\boxed{36}$.

 


 




	\item (\textit{AIME II 2015}) A cylindrical barrel with radius $4$ feet and height $10$ feet is full of water. A solid cube with side length $8$ feet is set into the barrel so that the diagonal of the cube is vertical. The volume of water thus displaced is $v$ cubic feet. Find $v^2$.


 
                    \begin{figure}[H]
                        \centering
                        \includegraphics[width=0.5\textwidth]{Problem/AIME_2015_II_9_Problem_Figure_1.png}
                    \end{figure}
                    


 
 \textbf{Answer: }384



\textbf{Solution 1}

Our aim is to find the volume of the part of the cube submerged in the cylinder. In the problem, since three edges emanate from each vertex, the boundary of the cylinder touches the cube at three points.  Because the space diagonal of the cube is vertical, by the symmetry of the cube, the three points form an equilateral triangle.  Because the radius of the circle is $4$, by the Law of Cosines, the side length s of the equilateral triangle is

 
 \[s^2 = 2\cdot(4^2) - 2\cdot(4^2)\cos(120^{\circ}) = 3(4^2)\]

 so $s = 4\sqrt{3}$.* Again by the symmetry of the cube, the volume we want to find is the volume of a tetrahedron with right angles on all faces at the submerged vertex, so since the lengths of the legs of the tetrahedron are $\frac{4\sqrt{3}}{\sqrt{2}} = 2\sqrt{6}$ (the three triangular faces touching the submerged vertex are all $45-45-90$ triangles) so 

 
 \[v = \frac{1}{3}(2\sqrt{6})\left(\frac{1}{2} \cdot (2\sqrt{6})^2\right) = \frac{1}{6} \cdot 48\sqrt{6} = 8\sqrt{6}\]

 so 

 
 \[v^2 = 64 \cdot 6 = \boxed{384}.\]

 In this case, our base was one of the isosceles triangles (not the larger equilateral one). To calculate volume using the latter, note that the height would be $2\sqrt{2}$.

 
                \begin{itemize}
                    \item Note that in a 30-30-120 triangle, side length ratios are $1:1:\sqrt{3}$.\item Or, note that the altitude and the centroid of an equilateral triangle are the same point, so since the centroid is 4 units from the vertex (which is $\frac{2}{3}$ the length of the median), the altitude is 6, which gives a hypotenuse of $\frac{12}{\sqrt{3}}=4\sqrt{3}$ by $1:\frac{\sqrt{3}}{2}:\frac{1}{2}$ relationship for 30-60-90 triangles.
                \end{itemize}


                





\textbf{Solution 2 (No trig)}

Visualizing the corner which is submerged in the cylinder, we can see that it's like slicing off a corner, where they slice an equal part off every edge, to make a tetrahedron where there are 3 right isosceles triangles and one equilateral triangle. With this out of the way, we can now just find the area of that equilateral triangle, using the fact that the circle of radius $4$ is the circumcircle of the equilateral triangle. Using equilateral triangle properties, you can find that the height of the triangle is $6$, and the side length is $\frac{6}{\sqrt{3}}=2\sqrt{3} \cdot 2=4\sqrt{3}$. As the other faces are right isosceles triangles, they are $\frac{4\sqrt{3}}{\sqrt{2}}=2\sqrt{6}$. Therefore the volume of this tetrahedron is

 
 \[\left(\frac{2\sqrt{6}}{2}\right)^2=12 \ \cdot \ (2\sqrt{6})=24\sqrt{6} \implies \frac{24\sqrt{6}}{3}=8\sqrt{6} \implies (8\sqrt{6})^2=\boxed{384}\]

 -dragoon

 -lpieleanu (minor latex changes)

 
Note: The height of the cylinder and the side length of the cube had no effect on the result of the problem.

 Note 2: If this doesn’t make sense, just imagine slicing a corner off a cube.

 





\textbf{Solution 3 (Also trig-less)}

We can use the same method as in Solution 2 to find the side length of the equilateral triangle, which is $4\sqrt3$. From here, its area is
 \[\dfrac{\bigl(4\sqrt3\bigr)^2\sqrt3}4=12\sqrt3.\]The leg of the isosceles right triangle is $\dfrac{4\sqrt3}{\sqrt2}=2\sqrt6$, and the horizontal distance from the vertex to the base of the tetrahedron is $4$ (the radius of the cylinder), so we can find the height, as shown in the diagram.
                    \begin{figure}[H]
                        \centering
                        \includegraphics[width=0.5\textwidth]{Solution/AIME_2015_II_9_Solution_Figure_1.png}
                    \end{figure}
                    The height from the tip to the base is $2\sqrt2$, so the volume is $\dfrac{12\sqrt3\cdot2\sqrt2}3=8\sqrt6$, and thus the answer is $\boxed{384}$.

 -integralarefun

 





\textbf{Solution 4 (De Gua's, no trig)}

Since the diagonal is perpendicular to the base of the cylinder, all three edges and faces can be treated symmetrically.

 The cross-section of the cube with the top face of the cylinder is an equilateral triangle inscribed in a circle with radius $4$. This means the medians of the triangle have length $\frac{3}{2} \cdot 4=6$, because the circumcenter is also the centroid, and the centroid divides the medians into lengths of ratio $2:1$. Using $30-60-90$ triangles, the side length of the triangle is $4\sqrt{3}$, and its area is $\frac{(4\sqrt{3})^2\sqrt{3}}{4}=12\sqrt{3}$. 

 Next, consider the submerged triangular sections of the faces. Each is a $45-45-90$ triangle with leg length $x$. The area of each is then $\frac{x^2}{2}$. By De Gua's Theorem on the submerged pyramid (which we can apply because it has a right-angled corner), $3\left( \frac{x^2}{2} \right) ^2=(12\sqrt{3})^2$. Solving yields $x=2\sqrt{6}$.

 The height of the pyramid is then $\sqrt{(2\sqrt{6})^2-4^2}=2\sqrt{2}$, by the Pythagorean Theorem (using the slant height and circumradius). The volume is then $v=\frac{1}{3}\cdot 12\sqrt{3} \cdot 2\sqrt{2}=8\sqrt{6}$, and the requested answer is $v^2=(8\sqrt{6})^2=\boxed{384}$. ~bad\_at\_mathcounts

 





\textbf{Solution 5(Similarity)}

First, note the diagonal is vertical so that the tetrahedron of cube in water is similar to the largest tetrahedron in the bottom of this cube. So we are trying to find their volume ratio by discovering their side length ration.

 In the equilateral triangle at the cylinder's top, we can find its side is 
 \[s_{1} = 2 * R * \sin(60^{\circ}) = 4\sqrt{3}\] as the base side of this tetrahedron.

 For the largest tetrahedron in the cube, its side length is just 
 \[s_{2} = 8\sqrt{2}\] a diagonal of one of the face. And its volume is 
 \[v_{2} = 8 * 8 * \frac{1}{2} * 8 * \frac{1}{3} = \frac{8^3}{6}\].

 Therefore, 
 \[v_{1} = v_{2} * (\frac{s_{1}}{s_{2}})^3 = 8\sqrt{6}\]
 \[v_{1}^2 = 384\]

 




	\item (\textit{AIME II 2015}) Call a permutation $a_1, a_2, \ldots, a_n$ of the integers $1, 2, \ldots, n$ \textit{quasi-increasing} if $a_k \leq a_{k+1} + 2$ for each $1 \leq k \leq n-1$. For example, $53421$ and $14253$ are quasi-increasing permutations of the integers $1, 2, 3, 4, 5$, but $45123$ is not. Find the number of quasi-increasing permutations of the integers $1, 2, \ldots, 7$.


 
 \textbf{Answer: }486



\textbf{Solution 1}

The simple recurrence can be found.

 When inserting an integer $n$ into a string with $n - 1$ integers, we notice that the integer $n$ has 3 spots where it can go: before $n - 1$, before $n - 2$, and at the very end.

 Ex. Inserting 4 into the string 123:4 can go before the 2 (1423), before the 3 (1243), and at the very end (1234).

 Only the addition of the next number, $n$, will change anything.

 Thus the number of permutations with $n$ elements is three times the number of permutations with $n-1$ elements. 

 Start with $n=3$ since all $6$ permutations work. And go up: $18, 54, 162, 486$.

 Thus for $n=7$ there are $2*3^5=\boxed{486}$ permutations.

 When you are faced with a brain-fazing equation and combinatorics is part of the problem, use recursion! This same idea appeared on another AIME with an 8-box problem.

 





\textbf{Solution 2 (Breif Solution of Solution 1)}

Now consider an integer $n$. We need to attach it to a permutation of length $n-1$. The integer $n$ should be placed before $n-1$ or $n-2$ or it should be completely placed at the end of the permutation of length $n-1$. Why? To satisfy the property that $a_k \leq a_{k+1} + 2$.

 Consider a permutation of length $2$. We have the possible permutations as $12, 21$.

 Now consider a permutation of length $3$. By the above mentioned strategy, the possible permutations are $312, 132, 123$ for the $12$ case and $231, 321, 213$ for the $21$ case.

 Observed It?

 Let $f(n)$ be the possible allowed permutations. Therefore $f(n) = 3 \cdot f(n-1)$.

 Now we have $f(2) = 2$.

 Hence using the relation we get:

 
 \[\begin{array}{|c|c|} \hline n & f(n) \\ \hline 2 & 2 \\ 3 & 6 \\ 4 & 18 \\ 5 & 54 \\ 6 & 162 \\ 7 & \boxed{486} \\ \hline \end{array}\]

 If you are interested in a direct relation then $f(n)=3^{n-2}\cdot 2$.

 $\sim$ Hardik (BlueE65)

 




	\item (\textit{AIME II 2015}) The circumcircle of acute $\triangle ABC$ has center $O$. The line passing through point $O$ perpendicular to $\overline{OB}$ intersects lines $AB$ and $BC$ at $P$ and $Q$, respectively. Also $AB=5$, $BC=4$, $BQ=4.5$, and $BP=\frac{m}{n}$, where $m$ and $n$ are relatively prime positive integers. Find $m+n$.


 
 \textbf{Answer: }023



\textbf{Solution 1}

Call $M$ and $N$ the feet of the altitudes from $O$ to $BC$ and $AB$, respectively. Let $OB = r$ . Notice that $\triangle{OMB} \sim \triangle{QOB}$ because both are right triangles, and $\angle{OBQ} \cong \angle{OBM}$. By $\frac{MB}{BO}=\frac{BO}{BQ}$, $MB = r\left(\frac{r}{4.5}\right) = \frac{r^2}{4.5}$. However, since $O$ is the circumcenter of triangle $ABC$, $OM$ is a perpendicular bisector by the definition of a circumcenter. Hence, $\frac{r^2}{4.5} = 2 \implies r = 3$. Since we know $BN=\frac{5}{2}$ and $\triangle BOP \sim \triangle BNO$, we have $\frac{BP}{3} = \frac{3}{\frac{5}{2}}$. Thus, $BP = \frac{18}{5}$.  $m + n=\boxed{023}$.

 





\textbf{Solution 2 (fastest)}

Minor arc $BC = 2A$ so $\angle{BOC}=2A$. Since $\triangle{BOC}$ is isosceles ($BO$ and $OC$ are radii), $\angle{CBO}=(180-2A)/2=90-A$. $\angle{CBO}=90-A$, so $\angle{BQO}=A$. From this we get that $\triangle{BPQ}\sim \triangle{BCA}$. So $\dfrac{BP}{BC}=\dfrac{BQ}{BA}$, plugging in the given values we get $\dfrac{BP}{4}=\dfrac{4.5}{5}$, so $BP=\dfrac{18}{5}$, and $m+n=\boxed{023}$.

 





\textbf{Solution 3}

Let $r=BO$. Drawing perpendiculars, $BM=MC=2$ and $BN=NA=2.5$. From there, 
 \[OM=\sqrt{r^2-4}\] Thus, 
 \[OQ=\frac{\sqrt{4r^2+9}}{2}\] Using $\triangle{BOQ}$, we get $r=3$. Now let's find $NP$. After some calculations with $\triangle{BON}$ ~ $\triangle{OPN}$, ${NP=11/10}$. Therefore, 
 \[BP=\frac{5}{2}+\frac{11}{10}=18/5\] $18+5=\boxed{023}$.

 





\textbf{Solution 4}

Let $\angle{BQO}=\alpha$. Extend $OB$ to touch the circumcircle at a point $K$. Then, note that $\angle{KAC}=\angle{CBK}=\angle{QBO}=90^\circ-\alpha$. But since $BK$ is a diameter, $\angle{KAB}=90^\circ$, implying $\angle{CAB}=\alpha$. It follows that $APCQ$ is a cyclic quadrilateral.

 Let $BP=x$. By Power of a Point, 
 \[5x=4\cdot\frac 9 2\implies x=\frac{18}{5}.\]The answer is $18+5=\boxed{023}$.

 





\textbf{Solution 5}

$\textit{Note: This is not a very good solution, but it is relatively natural and requires next to no thinking.}$

 Denote the circumradius of $ABC$ to be $R$, the circumcircle of $ABC$ to be $O$, and the shortest distance from $Q$ to circle $O$ to be $x$.  

 Using Power of a Point on $Q$ relative to circle $O$, we get that $x(x+2r) = 0.5 \cdot 4.5 = \frac{9}{4}$. Using Pythagorean Theorem on triangle $QOB$ to get $(x + r)^2 + r^2 = \frac{81}{4}$. Subtracting the first equation from the second, we get that $2r^2 = 18$ and therefore $r = 3$. Now, set $\cos{ABC} = y$. Using law of cosines on $ABC$ to find $AC$ in terms of $y$ and plugging that into the extended law of sines, we get $\frac{\sqrt{4^2 + 5^2 - 2 \cdot 4 \cdot 5 x}}{\sqrt{1 - x^2}} = 2R = 6$. Squaring both sides and cross multiplying, we get $36x^2 - 40x + 5 = 0$. Now, we get $x = \frac{10 \pm \sqrt{55}}{18}$ using quadratic formula. If you drew a decent diagram, $B$ is acute and therefore $x = \frac{10 + \sqrt{55}}{18}$(You can also try plugging in both in the end and seeing which gives a rational solution). Note that $BP = 3\frac{1}{\sin{OPB}} = \frac{3}{\cos{\angle ABC - \angle QBO}}.$ Using the cosine addition formula and then plugging in what we know about $QBO$, we get that $BP = \frac{162}{2\cos{B} + \sqrt{5}\sin{B}}$. Now, the hard part is to find what $\sin{B}$ is. We therefore want $\frac{\sqrt{324 - (10 + \sqrt{55})^2}}{18} = \frac{\sqrt{169 - 20\sqrt{55}}}{18}$. For the numerator, by inspection $(a + b\sqrt{55})^2$ will not work for integers $a$ and $b$. The other case is if there is $(a\sqrt{5} + b\sqrt{11})^2$. By inspection, $5\sqrt{5} - 2\sqrt{11}$ works. Therefore, plugging all this in yields the answer, $\frac{18}{5} \rightarrow \boxed{23}$. Solution by hyxue

 





\textbf{Solution 6}


                    \begin{figure}[H]
                        \centering
                        \includegraphics[width=0.5\textwidth]{Solution/AIME_2015_II_11_Solution_Figure_1.png}
                    \end{figure}
                    Reflect $A$, $P$ across $OB$ to points $A'$ and $P'$, respectively with $A'$ on the circle and $P, O, P'$ collinear. Now, $\angle A'CQ = 180^{\circ} - \angle A'CB = \angle A'AB = \angle P'PB$ by parallel lines. From here, $\angle P'PB = \angle PP'B = \angle A'P'Q$ as $P, P', Q$ collinear. From here, $A'P'QC$ is cyclic, and by power of a point we obtain $\frac{18}{5} \implies \boxed{023}$.~awang11's sol

 





\textbf{Solution 7 (Coordinate Geometry)}

Let $M$ be the midpoint of $AB$, and let $\ell$ be the perpendicular bisector of $\overline{AB}$. We establish a coordinate system with $B$ being the origin and $\overline{BC}$ being the $x$-axis. Here is our setup:\begin{align*}A&=(5\cos{B},5\sin{B})\\B&=(0,0)\\C&=(4,0)\\Q&=\left(\frac{9}{2},0\right)\\M&=\left(\frac{5}{2}\cos{B},\frac{5}{2}\sin{B}\right)\\O&=(x_{1},y_{1})\\P&=(x_{2},y_{2})\\AB&: y=\tan{B}x\\\ell&: y=-\frac{1}{\tan{B}}\left(x-\frac{5}{2}\cos{B}\right)+\frac{5}{2}\sin{B}\\\end{align*}
The circumcenter $O$ is the intersection of line $x=2$ and $\ell$. We find
 \[O=\left(2,\frac{5-4\cos{B}}{2\sin{B}}\right)\]Thus, we have the equation of line $PQ$:
 \[y=-\frac{4\sin{B}}{5-4\cos{B}}\left(x-\frac{9}{2}\right)\]This allows us to compute $x_{2}$:\begin{align*}\tan{(B)}\cdot x_{2}&=-\frac{4\sin{B}}{5-4\cos{B}}\left(x_{2}-\frac{9}{2}\right)\\x_{2}&=\frac{18\sin{B}}{5-4\cos{B}}\cdot\frac{1}{\left(\frac{\sin{B}}{\cos{B}}+\frac{4\sin{B}}{5-4\cos{B}}\right)}\\x_{2}&=\frac{18\sin{B}}{5-4\cos{B}}\cdot \frac{\cos{B}(5-4\cos{B})}{5\sin{B}}\\x_{2}&=\frac{18\cos{B}}{5}\end{align*}
We see that $x_2=BP\cdot \cos{B}$. Therefore, $BP=\frac{18}{5}$, and the requested sum is $18+5=\boxed{023}$.

 




	\item (\textit{AIME II 2015}) There are $2^{10} = 1024$ possible $10$-letter strings in which each letter is either an A or a B. Find the number of such strings that do not have more than $3$ adjacent letters that are identical.


 
 \textbf{Answer: }548



\textbf{Solution 1}

Let $f(n)$ be the number of valid arrangements in a string of length $n$.

 A string can end in 8 ways:

 
                \begin{itemize}
                    \item ...AAA\item ...BBB\item ...AAB\item ...BBA\item ...BAB\item ...ABA\item ...ABB\item ...BAA
                \end{itemize}


                
If it ends in ...AAA or ...BBB, the only condition is that the string of length $n-3$ preceding the 3 letters can't end with the same letters. If the string ends with ...AAA, the $n-3$ string must end with B in $f(n-3)/2$ ways, and if the string ends with ...BBB, the $n-3$ string must end with A in $f(n-3)/2$ ways, giving us a total of $f(n-3)$ ways.

 
If it ends in ...AAB or ...BBA, the $n-2$ string can't end with the 2 repeating letters coming after it. If it ends with ...AAB, the $n-2$ string can end in BB or AB or BA, and if it ends with ...BBA, the $n-2$ string can end in AA, AB, or BA. This happens in $4(f(n-2)/4) = f(n-2)$ ways.

 
Finally, if it ends in ...BAB, ...ABA, ...ABB, or ...BAA, the $n-1$ string must end in the last two letters (for example, if the string of length n ended in BAB, the n-1 string must end with BA). So this gives us $4(f(n-1)/4) = f(n-1)$ ways.

 
So $f(n) = f(n-1) + f(n-2) + f(n-3)$. We can calculate the base cases pretty easily: $f(1)$ is 2, $f(2)$ is 4, $f(3)$ is 8, so $f(4)$ is 14. Now we continue working upwards: $f(5) = 14 + 8 + 4 = 26$. $f(6) = 26 + 14 + 8 = 48$. $f(7) = 48 + 26 + 14 = 88$. $f(8) = 88 + 48 + 26 = 162$. $f(9) = 162 + 88 + 48 = 298$. Ans = $f(10) = 298 + 162 + 88 = \boxed{548}$.

 





\textbf{Solution 2}

This is a recursion problem. Let $a_n$ be the number of valid strings of $n$ letters, where the last letter is $A$.  Similarly, let $b_n$ be the number of valid strings of $n$ letters, where the last letter is $B$.

 Note that $a_n=b_{n-1}+b_{n-2}+b_{n-3}$ for all $n\ge4$.

 Similarly, we have $b_n=a_{n-1}+a_{n-2}+a_{n-3}$ for all $n\ge4$.

 Here is why:  every valid strings of $n$ letters $(n\ge4)$ where the first letter is $A$ must begin with one of the following:

 $AAAB$ - and the number of valid ways is $b_{n-3}$.

 $AAB$ - and the number of valid ways is $b_{n-2}$.

 $AB$ - and there are $b_{n-1}$ ways.

 We know that $a_1=1$, $a_2=2$, and $a_3=4$.  Similarly, we have $b_1=1$, $b_2=2$, and $b_3=4$.  We can quickly check our recursion to see if our recursive formula works.  By the formula, $a_4=b_3+b_2+b_1=7$, and listing out all $a_4$, we can quickly verify our formula.

 Therefore, we have the following:

 $\begin{tabular}{c|c|c|c|c|c|c|c|c|c|c}Value of n & 1 & 2 & 3 & 4 & 5 & 6 & 7 & 8 & 9 & 10\\\hline a & 1 & 2 & 4 & 7 & 13 & 24 & 44 & 81 & 149 & 274\\b & 1 & 2 & 4 & 7 & 13 & 24 & 44 & 81 & 149 & 274\end{tabular}$

 The total number of valid $10$ letter strings is equal to $a_{10}+b_{10}=274+274=\boxed{548}$.

 Notice that $a_n = b_n$, since $a_1=b_1$, $a_2=b_2$, and $a_3=b_3$.  Therefore, we didn't really need to list out both recursion formulas, which could save us some time.

 





\textbf{Solution 3 (Without Recursion)}

Playing around with strings gives this approach: We have a certain number of As, then Bs, and so on. Therefore, what if we denoted each solution with numbers like $(3,3,3,1)$ to denote AAABBBAAAB or vice versa (starting with Bs)? Every string can be represented like this!

 We can have from 4 to 10 numbers in our parentheses. For each case, we will start with the largest number possible, usually a bunch of 3s, then go down systematically. Realize also that if we are left with just 2s and 1s, there is only one number of 2s and 1s that adds up to the leftover amount. Our final answer is the sum of all of these parenthetical sets [each set multiplied by its permutations, as order matters] multiplied by two [starting with either A or B, and alternating as we go along].

 $4 \rightarrow (3,3,3,1) = 4, (3,3,2,2) = 6$

 $5 \rightarrow (3, 3, 2, 1, 1) = 30, (3, 2, 2, 2, 1) = 20, (2,2,2,2,2)=1$

 $6 \rightarrow (3,3,1,1,1,1) = 15, (3,2,2,1,1,1) = 60, (2,2,2,2,1,1) = 15$

 $7 \rightarrow (3,2,1,1,1,1,1) = 42, (2,2,2,1,1,1,1) = 35$

 $8 \rightarrow (3,1...1) = 8, (2,2,1...1) = 28$

 $9 \rightarrow (2,1...1) = 9$

 $10 \rightarrow (1,1....1) =1$

 Adding them all up gives you 274; multiplying by 2 gives $\boxed{548}$.

 





\textbf{Solution 4}

We are going to build the string, 1 character at a time. And, we are going to only care about the streak of letters at the end of the string.

 Let $a_n$ be the number of strings of length n that satisfy the problem statement and also has a "streak" of length 1 at the end. ABABBA has a streak of length 1.

 Let $b_n$ be the number of strings of length n that satisfy the problem statement and also has of length 2 at the end. ABABBAA has a streak of length 2. There are 2 "A" s at the end of the string.

 Let $c_n$ be the number of string of length n that that satisfy the problem statement and also has a "streak" of length 3 at the end. ABABBAAA has a streak of length 3. There are 3 "A" s at the end of the string.

 Let's establish a recursive relationship. $a_{n+1} = a_n+b_n+c_n$, since you can simply break the streak. $b_{n+1} = a_n$, and $c_{n+1} = b_n$ Since you can just add to the streak.

 We can log everything using a table.

 $\begin{tabular}{c|c|c|c|c|c|c|c|c|c|c}Value of n & 1 & 2 & 3 & 4 & 5 & 6 & 7 & 8 & 9 & 10\\\hline a & 2 & 2 & 4 & 8 & 14 & 26 & 48 & 88 & 162 & 298\\b & 0 & 2 & 2 & 4 & 8 & 14 & 26 & 48 & 88 & 162 \\c & 0 & 0 & 2 & 2 &4 & 8 & 14 & 26 & 48 & 88\end{tabular}$

 Adding $a_{10}$, $b_{10}$, $c_{10}$ gets the total number of numbers that doesn't have more than 3 concecutive letters.

 That gets a total of $298+162+88 = \boxed{548}$

 -AlexLikeMath

 





\textbf{Solution 5}

Let $S_n$ be the number of n-letter strings satisfying the problem criteria. Then we can easily see $S_1 = 2, S_2 = 4, S_3 = 8, S_4 = 16-2 =14$. For $n \ge 4$ consider the last three elements of the list. For example to get $S_5$ we could try to take $S_4$ and multiply by 2 since for each $S_4$ string we have two choices for the fifth character. However we have to be careful if the last three characters of $S_4$ are all the same, as in that case we would be overcounting, (i.e if we have a string $\dots aaa$, we can only add $b$ to the end so that the problem criteria is satisfied. Additionally strings such as $\dots aaa$ or $\dots bbb$ should only be counted once as we only have one choice for the $n$th character to add ($b$ and $a$ respectively). 

 Thus to compute $S_n$ we start with $S_{n-1}$ then take out the strings ending in $\dots bbb$ or $\dots aaa$. There are $S_{n-4}$ remaining valid strings (i.e if we pick any of the $S_{n-4}$ valid strings, examine the last character, if it is a $b$ then we append $aaa$ to it, and if if it an $a$ we append $bbb$ to it, hence the number of $S_{n-1}$ strings are in one-to-one correspondence with the number of strings in $S_{n-4}$. Thus we have the recursion $S_n = 2(S_{n-1} - S_{n-4}) + S_{n-4}$ (where we are first taking away $S_{n-4}$, doubling the result, then adding $S_{n-4}$ back in signifying that we only count it once, as described above). 

 The recursion simplifies to $S_n = 2S_{n-1} - S_{n-4}$ and we can now quickly compute the remaining values: 
 \[S_5 = 2S_4 - S_1 = 26, S_6 = 2(26)-4 = 48, S_7 = 2(48) - 8 = 88, S_8 = 2(88) - 14 = 162, S_9 = 2(162) - 26 = 298, S_{10} = 2S_9 - 48 = \boxed{548}\]

 





\textbf{Solution 6 (Generating Functions)}

Let us define a "run" as a set of consecutive letters that are all the same and let a "maximal run" be a run that is not the proper subset of any other run, in other words, it cannot be expanded.

 From now on, we consider all runs to be maximal.

 Note that the minimum number of runs in the $10$-letter string is $4$ (or else the Pigeonhole Principle tells us that at least $1$ run has $4$ or more letters in it, contradiction), and the maximum number of runs is clearly $10.$

 Since an arbitrary run can have anywhere from $1$ to $3$ letters in it, it follows that the number of $10$-letter strings with $n$ runs (where $4 \leq n \leq 10$) is
 \[2[x^{10}] (x+x^2+x^3)^n,\]i.e. twice the coefficient of $x^{10}$ in the expansion of $(x+x^2+x^3)^n.$ (Note that we multiplied by $2$ because there are two choices for which letter the first run has and then the rest are fixed).

 Hence, we wish to find
 \[2\sum_{n=4}^{10} [x^{10}] (x+x^2+x^3)^n.\]First, we can rewrite this as\begin{align*}\sum_{n=4}^{10} [x^{10}] (x+x^2+x^3)^n &= \sum_{n=4}^{10} [x^{10}] \ x^n \cdot (1+x+x^2)^n \\&= \sum_{n=4}^{10} [x^{10-n}] \left(\frac{1-x^3}{1-x}\right)^n \\&= \sum_{n=4}^{10} [x^{10-n}] \frac{(1-x^3)^n}{(1-x)^n}. \\\end{align*}
Now, we proceed case by case, utilizing the Binomial Theorem for the numerator in all of the cases:

 For $n=4,$ we have\begin{align*}[x^6] \frac{(1-x^3)^4}{(1-x)^4} &= \binom{4}{0} [x^6] (1-x)^{-4} - \binom{4}{1} [x^3] (1-x)^{-4} + \binom{4}{2} [x^0] (1-x)^{-4} \\&= 1\cdot\binom{9}{3} - 4\cdot\binom{6}{3} + 6\cdot\binom{3}{3} \\&= 1\cdot84-4\cdot20+6\cdot1 \\&= 84-80+6 \\&= \underline{10}.\end{align*}
For $n=5,$ we have\begin{align*}[x^5] \frac{(1-x^3)^5}{(1-x)^5} &= \binom{5}{0} [x^5] (1-x)^{-5} - \binom{5}{1}[x^2] (1-x)^{-5} \\&= 1\cdot\binom{9}{4}-5\cdot\binom{6}{4} \\&= 1\cdot126-5\cdot15 \\&= 126-75 \\&= \underline{51}.\end{align*}
For $n=6,$ we have\begin{align*}[x^4] \frac{(1-x^3)^6}{(1-x)^6} &= \binom{6}{0} [x^4](1-x)^{-6} - \binom{6}{1}[x^1](1-x)^{-6} \\&= 1 \cdot \binom{9}{5} - 6 \cdot \binom{6}{5} \\&= 1 \cdot 126 - 6 \cdot 6 \\&= 126-36 \\&= \underline{90}.\end{align*}
For $n=7,$ we have\begin{align*}[x^3] \frac{(1-x^3)^7}{(1-x)^7} &= \binom{7}{0}[x^3](1-x)^{-7} - \binom{7}{1}[x^0](1-x)^{-7} \\&= 1\cdot \binom{9}{6} - 7\cdot\binom{6}{6} \\&= 1\cdot84-7\cdot1\\&= 84-7 \\&= \underline{77}.\end{align*}
For $n=8,$ we have\begin{align*}[x^2] \frac{(1-x^3)^8}{(1-x)^8} &=\binom{8}{0} [x^2](1-x)^{-8} \\&= 1\cdot\binom{9}{7} \\&= 1\cdot36 \\&= \underline{36}.\end{align*}
For $n=9,$ we have\begin{align*}[x^1] \frac{(1-x^3)^9}{(1-x)^9} &= \binom{9}{0} [x^1] (1-x)^{-9}\\&= 1\cdot\binom{9}{8} \\&= 1\cdot9 \\&= \underline{9}.\end{align*}
For $n=10,$ we have\begin{align*}[x^0] \frac{(1-x^3)^{10}}{(1-x)^{10}} &= [x^0] (1-x)^{-10} \\&= 1\cdot\binom{9}{9} \\&= 1\cdot1 \\&= \underline{1}.\end{align*}
Hence, the answer is\begin{align*}2\sum_{n=4}^{10} [x^{10-n}]  \frac{(1-x^3)^n}{(1-x)^n} &= 2(10+51+90+77+36+9+1) \\&= 2\cdot274 \\&= \boxed{548}.\end{align*}


 -lpieleanu

 \textbf{Remark} 2015 AMC 12A \#22 is a harder version of this problem.

 




	\item (\textit{AIME II 2015}) Define the sequence $a_1, a_2, a_3, \ldots$ by $a_n = \sum\limits_{k=1}^n \sin{k}$, where $k$ represents radian measure. Find the index of the 100th term for which $a_n < 0$.


 
 \textbf{Answer: }628



\textbf{Solution 1}

If $n = 1$, $a_n = \sin(1) > 0$.  Then if $n$ satisfies $a_n < 0$, $n \ge 2$, and
 \[a_n = \sum_{k=1}^n \sin(k) = \cfrac{1}{\sin{1}} \sum_{k=1}^n\sin(1)\sin(k) = \cfrac{1}{2\sin{1}} \sum_{k=1}^n\cos(k - 1) - \cos(k + 1) = \cfrac{1}{2\sin(1)} [\cos(0) + \cos(1) - \cos(n) - \cos(n + 1)].\]Since $2\sin 1$ is positive, it does not affect the sign of $a_n$.  Let $b_n = \cos(0) + \cos(1) - \cos(n) - \cos(n + 1)$.  Now since $\cos(0) + \cos(1) = 2\cos\left(\cfrac{1}{2}\right)\cos\left(\cfrac{1}{2}\right)$ and $\cos(n) + \cos(n + 1) = 2\cos\left(n + \cfrac{1}{2}\right)\cos\left(\cfrac{1}{2}\right)$, $b_n$ is negative if and only if $\cos\left(\cfrac{1}{2}\right) < \cos\left(n + \cfrac{1}{2}\right)$, or when $n \in [2k\pi - 1, 2k\pi]$.  Since $\pi$ is irrational, there is always only one integer in the range, so there are values of $n$ such that $a_n < 0$ at $2\pi, 4\pi, \cdots$.  Then the hundredth such value will be when $k = 100$ and $n = \lfloor 200\pi \rfloor = \lfloor 628.318 \rfloor = \boxed{628}$.

 \textbf{Clearer Explanation for the End} Notice that we want $\cos(n+1/2)>\cos(1/2).$

 Notice that the only interval in which this is true, is where $-1/2 < n + 1/2 < 1/2.$

 Thus, we need $-1<n<0.$ So we have $-1+2\pi k <n <0+2\pi k.$

 But we want the 100th solution. Which means $n$ is in between $200\pi - 1 < n < 200 \pi.$ So n = 628.

 





\textbf{Solution 2}

Notice that $a_n$ is the imaginary part of $\sum_{k=1}^n e^{ik}$, by Euler's formula. Using the geometric series formula, we find that this sum is equal to 
 \[\frac{e^{i(n+1)}-e^i}{e^i-1} = \frac{\cos (n+1) - \cos 1 + i (\sin (n+1) - \sin 1) }{\cos 1 - 1 + i \sin 1}\] We multiply the fraction by the conjugate of the denominator so that we can separate out the real and imaginary parts of the above expression. Multiplying, we have 
 \[\frac{(\cos 1 - 1)(\cos(n+1)-\cos 1) + (\sin 1)(\sin(n+1)-\sin 1) + i((\sin(n+1) - \sin 1)(\cos 1 - 1) - (\sin 1)(\cos(n+1)-\cos 1))}{\cos^2 1 - 2 \cos 1 + 1 + \sin^2 1}\] We only need to look at the imaginary part, which is 
 \[\frac{(\sin(n+1) \cos 1 - \cos(n+1) \sin 1) - \sin 1 \cos 1 + \sin 1 - \sin (n+1) + \sin 1 \cos 1}{2-2 \cos 1} = \frac{\sin n - \sin(n+1) + \sin 1}{2-2 \cos 1}\]Since $\cos 1 < 1$, $2-2 \cos 1 > 0$, so the denominator is positive. Thus, in order for the whole fraction to be negative, we must have $\sin (n+1) - \sin n > \sin 1 \implies 2 \cos \left(n + \frac{1}{2} \right) \sin \frac{1}{2} > \sin 1 \implies \cos \left( n + \frac{1}{2} \right) > \frac{\sin 1}{2 \sin{\frac{1}{2}}} = \cos \left(\frac{1}{2} \right),$ by sum to product. This only holds when $n$ is between $2\pi k - 1$ and $2\pi k$ for integer $k$ [continuity proof here], and since this has exactly one integer solution for every such interval, the $100$th such $n$ is $\lfloor 200\pi \rfloor = \boxed{628}$.

 





\textbf{Solution 3}

Similar to solution 2, we set a complex number $z=\cos 1+i\sin 1$. We start from $z$ instead of $1$ because $k$ starts from $1$: be careful.

 The sum of $z+z^2+z^3+z^4+z^5\dots=\frac{z-z^{n+1}}{1-z}=\frac{z}{z-1}\left(z^n-1\right)$.

 We are trying to make $n$ so that the imaginary part of this expression is negative.

 The argument of $z$ is $1$. The argument of $z-1$, however, is a little more tricky. $z-1$ is on a circle centered on $(-1,0)$ with radius $1$. The change in angle due to $z$ is $1$ with respect to the center, but the angle that $z-1$ makes with the $y$-axis is $half$ the change, due to Circle Theorems (this intercepted arc is the argument of $z$), because the $y$- axis is tangent to the circle at the origin. So $\text{arg}(z-1)=\frac{\pi+1}{2}$. Dividing $z$ by $z-1$ subtracts the latter argument from the former, so the angle of the quotient with the $x$-axis is $\frac{1-\pi}{2}$. 

 We want the argument of the whole expression $-\pi<\theta<0$. This translates into $\frac{-\pi-1}{2}<\text{arg}\left(z^n-1\right)<\frac{\pi-1}{2}$. $z^n-1$ also consists of points on the circle centered at $(-1,0)$, so we deal with this argument similarly: the argument of $z^n$ is twice the angle $z^n-1$ makes with the $y$-axis. Since $z^n-1$ is always negative, $\frac{-3\pi}{2}<\text{arg}\left(z^n-1\right)<\frac{-\pi}{2}$, and the left bound is the only one that is important. Either way, the line (the line consists of both bounds) makes a $\frac{\pi}{2}-\frac{\pi-1}{2}=\frac{-1}{2}$ angle with the $y$-axis both ways.

 So the argument of $z^n$ must be in the bound $-1<\theta<0$ by doubling, namely the last $z^n$ negative before another rotation. Since there is always one $z^n$ in this category per rotation because $\pi$ is irrational, $n_{100}\equiv z^{628}$ and the answer is $\boxed{628}$.

 





\textbf{Solution 4}

By the product-to-sum formula,
 \[\sin \frac{1}{2} \sin k = \frac{1}{2} \left[ \cos \left( k - \frac{1}{2} \right) - \cos \left( k + \frac{1}{2} \right) \right].\]Thus, we can make the sum in the problem telescope:\begin{align*}a_n &= \sum_{k = 1}^n \sin k \\&= \sum_{k = 1}^n \frac{\sin \frac{1}{2} \sin k}{\sin \frac{1}{2}} \\&= \sum_{k = 1}^n \frac{\cos (k - \frac{1}{2}) - \cos (k + \frac{1}{2})}{2 \sin \frac{1}{2}} \\&= \frac{(\cos \frac{1}{2} - \cos \frac{3}{2}) + (\cos \frac{3}{2} - \cos \frac{5}{2}) + \dots + (\cos \frac{2n - 1}{2} - \cos \frac{2n + 1}{2})}{2 \sin \frac{1}{2}} \\&= \frac{\cos \frac{1}{2} - \cos \frac{2n + 1}{2}}{2 \sin \frac{1}{2}}.\end{align*}
Then $a_n < 0$ when $\cos \frac{1}{2} < \cos \frac{2n + 1}{2}.$ This occurs if and only if
 \[2 \pi k - \frac{1}{2} < \frac{2n + 1}{2} < 2 \pi k + \frac{1}{2}\]for some integer $k.$ Equivalently,
 \[2 \pi k - 1 < n < 2 \pi k.\]In other words, $n = \lfloor 2 \pi k \rfloor.$ The 100th index of this form is then $\lfloor 2 \pi \cdot 100 \rfloor = \boxed{628}.$

 




	\item (\textit{AIME II 2015}) Let $x$ and $y$ be real numbers satisfying $x^4y^5+y^4x^5=810$ and $x^3y^6+y^3x^6=945$. Evaluate $2x^3+(xy)^3+2y^3$.


 
 \textbf{Answer: }089



\textbf{Solution 1}

The expression we want to find is $2(x^3+y^3) + x^3y^3$.

 Factor the given equations as $x^4y^4(x+y) = 810$ and $x^3y^3(x^3+y^3)=945$, respectively. Dividing the latter by the former equation yields $\frac{x^2-xy+y^2}{xy} = \frac{945}{810}$. Adding 3 to both sides and simplifying yields $\frac{(x+y)^2}{xy} = \frac{25}{6}$. Solving for $x+y$ and substituting this expression into the first equation yields $\frac{5\sqrt{6}}{6}(xy)^{\frac{9}{2}} = 810$. Solving for $xy$, we find that $xy = 3\sqrt[3]{2}$, so $x^3y^3 = 54$. Substituting this into the second equation and solving for $x^3+y^3$ yields $x^3+y^3=\frac{35}{2}$. So, the expression to evaluate is equal to $2 \times \frac{35}{2} + 54 = \boxed{089}$.

 Note that since the value we want to find is $2(x^3+y^3)+x^3y^3$, we can convert $2(x^3+y^3)$ into an expression in terms of $x^3y^3$, since from the second equation which is $x^3y^3(x^3+y^3)=945$, we see that $2(x^3+y^3)=1890+x^6y^6,$ and thus the value is $\frac{1890+x^6y^6}{x^3y^3}.$ Since we've already found $x^3y^3,$ we substitute and find the answer to be 89.

 





\textbf{Solution 2}

Factor the given equations as $x^4y^4(x+y) = 810$ and $x^3y^3(x+y)(x^2-xy+y^2)=945$, respectively. By the first equation, $x+y=\frac{810}{x^4y^4}$. Plugging this in to the second equation and simplifying yields $(\frac{x}{y}-1+\frac{y}{x})=\frac{7}{6}$. Now substitute $\frac{x}{y}=a$. Solving the quadratic in $a$, we get $a=\frac{x}{y}=\frac{2}{3}$ or $\frac{3}{2}$ As both of the original equations were symmetric in $x$ and $y$, WLOG, let $\frac{x}{y}=\frac{2}{3}$, so $x=\frac{2}{3}y$. Now plugging this in to either one of the equations, we get the solutions $y=\frac{3(2^{\frac{2}{3}})}{2}$, $x=2^{\frac{2}{3}}$. Now plugging into what we want, we get $8+54+27=\boxed{089}$

 





\textbf{Solution 3}

Add three times the first equation to the second equation and factor to get $(xy)^3(x^3+3x^2y+3xy^2+y^3)=(xy)^3(x+y)^3=3375$. Taking the cube root yields $xy(x+y)=15$. Noting that the first equation is $(xy)^3\cdot(xy(x+y))=810$, we find that $(xy)^3=\frac{810}{15}=54$. Plugging this into the second equation and dividing yields $x^3+y^3 = \frac{945}{54} = \frac{35}{2}$. Thus the sum required, as noted in Solution 1, is $54+\frac{35}{2}\cdot2 = \boxed{089}$.

 





\textbf{Solution 4}

As with the other solutions, factor. But this time, let $a=xy$ and $b=x+y$. Then $a^4b=810$. Notice that $x^3+y^3 = (x+y)(x^2-xy+y^2) = b(b^2-3a)$, so we are looking for $2b(b^2-3a)+a^3$. Now, if we divide the second equation by the first one, we get $7/6 = \frac{b^2-3a}{a}$; then $\frac{b^2}{a}=\frac{25}{6}$. Therefore, $a = \frac{6}{25}b^2$. Substituting $\frac{6}{25}b^2$ for $a$ in equation 1, simplifying, and then taking the cube root gives us $b^3 = \frac{5^3}{2}.$ Finding $a^3$ by cubing $a = \frac{6}{25}b^2$ on both sides and simplifying using our previous substitution, we get $a^3 = 54$. Substituting this into the first equation and then dividing by $27$, we get $2b(b^2 - 3a) = 35$. Our final answer is $35+54=\boxed{089}$.

 





\textbf{Solution 5}

Factor the given equations as:
 \[x^4y^4(x+y)=810\]
 \[x^3y^3(x^3+y^3)=x^3y^3(x+y)(x^2-xy+y^2)=945\]We note that these expressions (as well as the desired expression) can be written exclusively in terms of $x+y$ and $xy$. We make the substitution $s=x+y$ and $p=xy$ (for sum and product, respectively).

 
 \[x^4y^4(x+y)=p^4s=810\]
 \[x^3y^3(x+y)(x^2-xy+y^2)=p^3(s)(s^2-3p)=s^3p^3-3p^4s=945\]

 We see that $p^4s$ shows up in both equations, so we can eliminate it and find $sp$, after which we can get $p^3$ from the first equation. If you rewrite the desired expression using $s$ and $p$, it becomes clear that you don't need to actually find the values of $s$ and $p$, but I will do so for the sake of completion.

 
 \[s^3p^3=945+3p^4s\]
 \[s^3p^3=945+3(810)=3375\]
 \[sp=15\]

 
 \[p^3=\frac{810}{sp}=54\]
 \[p=3\cdot2^{1/3}\]
 \[s=\frac{15}{p}=5\cdot2^{-1/3}\]

 The desired expression can be written as:
 \[2(x^3+y^3)+(xy)^3=2(x+y)(x^2-xy+y^2)+(xy)^3\]
 \[2(s)(s^2-3p)+p^3=2s^3-6sp+p^3\]

 Plugging in $s$ and $p$, we get:
 \[2(5\cdot2^{-1/3})^3-6(15)+54=125-90+54=\boxed{089}\]

 - gting

 





\textbf{Solution 6}

Factor the first and second equations as $(xy)^4(x+y)=810$ and $(xy)^3(x+y)(x^2-xy+y^2)=945$, respectively. Dividing them (allowed, since neither are $0$), we have
 \[\frac{xy}{x^2-xy+y^2}=\frac67\]or
 \[x^2-\frac{13}{6}xy+y^2=0.\]Plugging into the quadratic formula and solving for $x$ in terms of $y,$ we have
 \[x=\frac{\frac{13y}{6}\pm \sqrt{\frac{169y^2}{36}-4y^2}}{2}=\frac{2y}3 , \frac{3y}2 .\]WLOG, let $x=\frac{3y}2 .$Plugging into our first equation, we have
 \[\left(\frac{3}{2}y\right)^4\left(\frac52 y\right)=810 \implies y^3 = 4 .\]Plugging this result (and the one for $x$ in terms of $y$) into our desired expression, we have

 





\textbf{Solution 7}

Take $w=x+y$ and $z=xy$. Remark that 
 \begin{align*} ~&2x^{3}+(xy)^{3}+2y^{3} \\ =~&2(x^{3}+y^{3})+(xy)^{3} \\ =~&2\left[(x+y)^{3}-3xy(x+y)\right]+(xy)^{3} \\ =~&2(w^{3}-3wz)+z^{3} \\ =~& 2w^{3}-6wz+z^{3}.\end{align*}
 The given equations imply that 
 \[wz^{4}=810~~~\text{and}~~~(wz)^{3}-3wz^{4}=945.\] Substituting the first equation into the second, we have that $(wz)^{3}=945+3\cdot 810=3375$, thus $wz=\sqrt[3]{3375}=15$. Now 
 \[z^{3}=\frac{wz^{4}}{wz}=\frac{810}{15}=54\] and 
 \[w^{3}=\frac{3375}{54}=\frac{125}{2}.\] Thus 
 \begin{align*}2w^{3}-6wz+z^{3}&=2\left(\frac{125}{2}\right)-6(15)+54 \\ &=125-90+54 \\ &=\boxed{89}.\end{align*}


 


 





\textbf{Solution 8}

$x^4y^4(x+y)=810; x^3y^3(x^3+y^3)=945, \frac{x^2-xy+y^2}{xy}=\frac{7}{6}, \frac{x^2+y^2}{xy}=\frac{13}{6}$

 Let $x^2+y^2=13k; xy=6k$, then we can see $(x+y)^2-12k=13k, x+y=5\sqrt{k}$, now, we see $x^4y^4\cdot (x+y)=1296k^4\cdot 5\sqrt{k}, k=\frac{1}{\sqrt[3]{4}}$.

 The rest is easy, $2(x^3+y^3)+x^3y^3=216k^3+2[(x+y)^3-3xy(x+y)]=216k^3+2\cdot 35k^{\frac{3}{2}}=\boxed{89}$

 





\textbf{Solution 9}

Let's first put the left hand sides of the equations in factored forms. Doing this we obtain $(xy)^4 \cdot (x + y) = 810$ and $(xy)(x + y)(x^2 - xy + y^2) = 945$. Now, we will subtract and add the equations to gather information on $x$ and $y$. When we subtract the equations and clean it up via factoring, we yield $(xy)^3 \cdot (x + y) \cdot (x - y)^2 = 115$, and when we add them, we yield $(xy)^3 \cdot (x + y) \cdot (x^2 + y^2) = 1755$. Now with some intuition, you should divide the equations to obtain $\frac{(x^2 + y^2)}{(x - y)^2} = 13$. Now, we clean this up to obtain the following factoring of $0 = 2 \cdot (2x - 3y) \cdot (3x - 2y)$. This implies that $x = \frac{3y}{2}$. We plug that into the target expression to reduce it down to one variable, and get that target expression is $2x^3 + (xy)^3 + 2y^3 = \frac{27}{4} \cdot y^3 + \frac{27}{8} \cdot y^6 + 2y^3$. This means that if we can find a way to get $y^3$, then the rest is trivial. We get $y^3$ by plugging in $x = \frac{3y}{2}$ into $x^3 \cdot y^6 + y^3 \cdot x^6 = 945$. However, this time we only factor as $(xy)^3 \cdot (x^3 + y^3)$ because we particularly want a cubic degree on $y$. Plugging in $x = \frac{3y}{2}$ we get $y^3 = 4$. Now lets plug this into our target expression to get $\frac{27}{4} \cdot 4 + \frac{27}{8} \cdot 4^2 + 2 \cdot 4 = \boxed{89}$.

 





\textbf{Solution 10 (Newton sums)}

Let $x, y$ be roots of some polynomial $P(\alpha)$. Then we have that $P(\alpha) = (\alpha - x)(\alpha - y) \implies P(\alpha) = \alpha^2 - (x + y) \alpha + xy$. Now from the first equation we have that $x + y = 810/(xy)^4$. For convenience, denote $z = xy$.

 Define $P_k = x^k + y^k$, where $x, y$ are roots of the above polynomial. Then we have that by Newton's formulas $P_k = (x + y)P_{k - 1} - xyP_{k - 2} = 810/z^4P_{k - 1} - zP_{k - 2}$, where $P_0 = 2$ and $P_1 = x + y = 810/z^4$. We desire $2P_3 + z^3$. Now building up this recurrence; 

 
 \[P_2 = \frac{810}{z^4}P_1 - zP_0 = \frac{810^2}{z^8} - 2z.\] Then we have that 
 \[P_3 = \frac{810}{z^4} \left(\frac{810^2}{z^8} - 3z \right).\] Plugging this value of $P_3$ into the second equation $z^3P_3 = 810$, yields $z^9 = 164025$. Now $z^3 = 54$. To compute $P_3$ is now trivial, as returning this new value of $z$ into the second equation yields $P_3 = 35/2$. Hence $z^3 + 2P_3 = 54 + 35 = \boxed{089}$.

 -th1nq3r

 




	\item (\textit{AIME II 2015}) Circles $\mathcal{P}$ and $\mathcal{Q}$ have radii $1$ and $4$, respectively, and are externally tangent at point $A$. Point $B$ is on $\mathcal{P}$ and point $C$ is on $\mathcal{Q}$ such that $BC$ is a common external tangent of the two circles. A line $\ell$ through $A$ intersects $\mathcal{P}$ again at $D$ and intersects $\mathcal{Q}$ again at $E$. Points $B$ and $C$ lie on the same side of $\ell$, and the areas of $\triangle DBA$ and $\triangle ACE$ are equal. This common area is $\frac{m}{n}$, where $m$ and $n$ are relatively prime positive integers. Find $m+n$.


 
                    \begin{figure}[H]
                        \centering
                        \includegraphics[width=0.5\textwidth]{Problem/AIME_2015_II_15_Problem_Figure_1.png}
                    \end{figure}
                    


 
 \textbf{Answer: }129



\textbf{Solution 1 (guys trig is fast)}

Let $M$ be the intersection of $\overline{BC}$ and the common internal tangent of $\mathcal P$ and $\mathcal Q.$ We claim that $M$ is the circumcenter of right $\triangle{ABC}.$ Indeed, we have $AM = BM$ and $BM = CM$ by equal tangents to circles, and since $BM = CM, M$ is the midpoint of $\overline{BC},$ implying that $\angle{BAC} = 90.$ Now draw $\overline{PA}, \overline{PB}, \overline{PM},$ where $P$ is the center of circle $\mathcal P.$ Quadrilateral $PAMB$ is cyclic, and by Pythagorean Theorem $PM = \sqrt{5},$ so by Ptolemy on $PAMB$ we have 
 \[AB \sqrt{5} = 2 \cdot 1 + 2 \cdot 1 = 4 \iff AB = \dfrac{4 \sqrt{5}}{5}.\] Do the same thing on cyclic quadrilateral $QAMC$ (where $Q$ is the center of circle $\mathcal Q$ and get $AC = \frac{8 \sqrt{5}}{5}.$ 

 Let $\angle A = \angle{DAB}.$ By Law of Sines, $BD = 2R \sin A = 2 \sin A.$ Note that $\angle{D} = \angle{ABC}$ from inscribed angles, so 
 \begin{align*} [ABD] &= \dfrac{1}{2} BD \cdot AB \cdot \sin{\angle B} \\ &= \dfrac{1}{2} \cdot \dfrac{4 \sqrt{5}}{5} \cdot 2 \sin A \sin{\left(180 - \angle A - \angle D\right)} \\ &= \dfrac{4 \sqrt{5}}{5} \cdot \sin A \cdot \sin{\left(\angle A + \angle D\right)} \\ &= \dfrac{4 \sqrt{5}}{5} \cdot \sin A \cdot \left(\sin A \cos D + \cos A \sin D\right) \\ &= \dfrac{4 \sqrt{5}}{5} \cdot \sin A \cdot \left(\sin A \cos{\angle{ABC}} + \cos A \sin{\angle{ABC}}\right) \\ &= \dfrac{4 \sqrt{5}}{5} \cdot \sin A \cdot \left(\dfrac{\sqrt{5} \sin A}{5} + \dfrac{2 \sqrt{5} \cos A}{5}\right) \\ &= \dfrac{4}{5} \cdot \sin A \left(\sin A + 2 \cos A\right) \end{align*}
 after angle addition identity.

 Similarly, $\angle{EAC} = 90 - \angle A,$ and by Law of Sines $CE = 8 \sin{\angle{EAC}} = 8 \cos A.$ Note that $\angle{E} = \angle{ACB}$ from inscribed angles, so
 \begin{align*} [ACE] &= \dfrac{1}{2} AC \cdot CE \sin{\angle C} \\ &= \dfrac{1}{2} \cdot \dfrac{8 \sqrt{5}}{5} \cdot 8 \cos A \sin{\left[180 - \left(90 - \angle A\right) - \angle E\right]} \\ &= \dfrac{32 \sqrt{5}}{5} \cdot \cos A \sin{\left[\left(90 - \angle A\right) + \angle{ACB}\right]} \\ &= \dfrac{32 \sqrt{5}}{5} \cdot \cos A \left(\dfrac{2 \sqrt{5} \cos A}{5} + \dfrac{\sqrt{5} \sin A}{5}\right) \\ &= \dfrac{32}{5} \cdot \cos A \left(\sin A + 2 \cos A\right) \end{align*}
 after angle addition identity.Setting the two areas equal, we get 
 \[\tan A = \frac{\sin A}{\cos A} = 8 \iff \sin A = \frac{8}{\sqrt{65}}, \cos A = \frac{1}{\sqrt{65}}\] after Pythagorean Identity. Now plug back in and the common area is $\frac{64}{65} \iff \boxed{129}.$

 





\textbf{Solution 2}


                    \begin{figure}[H]
                        \centering
                        \includegraphics[width=0.5\textwidth]{Solution/AIME_2015_II_15_Solution_Figure_1.png}
                    \end{figure}
                    

 Call $O_1$ and $O_2$ the centers of circles $\mathcal{P}$ and  $\mathcal{Q}$, respectively, and extend $CB$ and $O_2O_1$ to meet at point $N$. Call $K$ and $L$ the feet of the altitudes from $B$ to $O_1N$ and $C$ to $O_2N$, respectively. Using the fact that $\triangle{O_1BN} \sim \triangle{O_2CN}$ and setting $NO_1 = k$, we have that $\frac{k+5}{k} = \frac{4}{1} \implies k=\frac{5}{3}$. We can do some more length chasing using triangles similar to $O_1BN$ to get that $AK = AL = \frac{24}{15}$, $BK = \frac{12}{15}$, and $CL = \frac{48}{15}$. Now, consider the circles $\mathcal{P}$ and $\mathcal{Q}$ on the coordinate plane, where $A$ is the origin. If the line $\ell$ through $A$ intersects $\mathcal{P}$ at $D$ and $\mathcal{Q}$ at $E$ then $4 \cdot DA = AE$. To verify this, notice that $\triangle{AO_1D} \sim \triangle{EO_2A}$ from the fact that both triangles are isosceles with $\angle{O_1AD} \cong \angle{O_2AE}$, which are corresponding angles. Since $O_2A = 4\cdot O_1A$, we can conclude that $4 \cdot DA = AE$.

 
Hence, we need to find the slope $m$ of line $\ell$ such that the perpendicular distance $n$ from $B$ to $AD$ is four times the perpendicular distance $p$ from $C$ to $AE$. This will mean that the product of the bases and heights of triangles $ACE$ and $DBA$ will be equal, which in turn means that their areas will be equal. Let the line $\ell$ have the equation $y = -mx \implies mx + y = 0$, and let $m$ be a positive real number so that the negative slope of $\ell$ is preserved. Setting $A = (0,0)$, the coordinates of $B$ are $(x_B, y_B) = \left(\frac{-24}{15}, \frac{-12}{15}\right)$, and the coordinates of $C$ are $(x_C, y_C) = \left(\frac{24}{15}, \frac{-48}{15}\right)$. Using the point-to-line distance formula and the condition $n = 4p$, we have 
 \[\frac{|mx_B + 1(y_B) + 0|}{\sqrt{m^2 + 1}} = \frac{4|mx_C + 1(y_C) + 0|}{\sqrt{m^2 + 1}}\] 
 \[\implies |mx_B + y_B| = 4|mx_C + y_C| \implies \left|\frac{-24m}{15} + \frac{-12}{15}\right| = 4\left|\frac{24m}{15} + \frac{-48}{15}\right|.\] If $m > 2$, then clearly $B$ and $C$ would not lie on the same side of $\ell$. Thus since $m > 0$, we must switch the signs of all terms in this equation when we get rid of the absolute value signs. We then have 
 \[\frac{6m}{15} + \frac{3}{15} = \frac{48}{15} - \frac{24m}{15} \implies 2m = 3 \implies m = \frac{3}{2}.\] Thus, the equation of $\ell$ is $y = -\frac{3}{2}x$.

 
Then we can find the coordinates of $D$ by finding the point $(x,y)$ other than $A = (0,0)$ where the circle $\mathcal{P}$ intersects $\ell$. $\mathcal{P}$ can be represented with the equation $(x + 1)^2 + y^2 = 1$, and substituting $y = -\frac{3}{2}x$ into this equation yields $x = 0, -\frac{8}{13}$ as solutions. Discarding $x = 0$, the $y$-coordinate of $D$ is  $-\frac{3}{2} \cdot -\frac{8}{13} = \frac{12}{13}$. The distance from $D$ to $A$ is then $\frac{4}{\sqrt{13}}.$ The perpendicular distance from $B$ to $AD$ or the height of $\triangle{DBA}$ is $\frac{|\frac{3}{2}\cdot\frac{-24}{15} + \frac{-12}{15} + 0|}{\sqrt{\frac{3}{2}^2 + 1}} = \frac{\frac{48}{15}}{\frac{\sqrt{13}}{2}} = \frac{32}{5\sqrt{13}}.$ Finally, the common area is $\frac{1}{2}\left(\frac{32}{5\sqrt{13}} \cdot \frac{4}{\sqrt{13}}\right) = \frac{64}{65}$, and $m + n = 64 + 65 = \boxed{129}$.

 





\textbf{Solution 3}

By \href{https://artofproblemsolving.com/wiki/index.php?title=Homothety}{homothety}, we deduce that $AE = 4 AD$. (The proof can also be executed by similar triangles formed from dropping perpendiculars from the centers of $P$ and $Q$ to $l$.) Therefore, our equality of area condition, or the equality of base times height condition, reduces to the fact that the distance from $B$ to $l$ is four times that from $C$ to $l$. Let the distance from $C$ be $x$ and the distance from $B$ be $4x$.

 Let $P$ and $Q$ be the centers of their respective circles. Then dropping a perpendicular from $P$ to $Q$ creates a $3-4-5$ right triangle, from which $BC = 4$ and, if $\alpha = \angle{AQC}$, that $\cos \alpha = \dfrac{3}{5}$. Then $\angle{BPA} = 180^\circ - \alpha$, and the Law of Cosines on triangles $APB$ and $AQC$ gives $AB = \dfrac{4}{\sqrt{5}}$ and $AC = \dfrac{8}{\sqrt{5}}.$

 Now, using the Pythagorean Theorem to express the length of the projection of $BC$ onto line $l$ gives
 \[\sqrt{\frac{16}{5} - 16x^2} + \sqrt{\frac{64}{5} - x^2} = \sqrt{16 - 9x^2}.\]Squaring and simplifying gives
 \[\sqrt{\left(\frac{1}{5} - x^2\right)\left(\frac{64}{5} - x^2\right)} = x^2,\]and squaring and solving gives $x = \dfrac{8}{5\sqrt{13}}.$

 By the Law of Sines on triangle $ABD$, we have
 \[\frac{BD}{\sin A} = 2.\]But we know $\sin A = \dfrac{4x}{AB}$, and so a small computation gives $BD = \dfrac{16}{\sqrt{65}}.$ The Pythagorean Theorem now gives
 \[AD = \sqrt{BD^2 - (4x)^2} + \sqrt{AB^2 - (4x)^2} = \frac{4}{\sqrt{13}},\]and so the common area is $\dfrac{1}{2} \cdot \frac{4}{\sqrt{13}} \cdot \frac{32}{5\sqrt{13}} = \frac{64}{65}.$ The answer is $\boxed{129}.$

 \textbf{Alternate Path to x} Call the intersection of lines $l$ and $BC$ $E$.You can use similar triangles to find that the distance from $B$ to $E$ is four times the distance from $C$ to $E$. Then draw a perpendicular from $A$ to $BC$ and call the point $F$. $AF = \frac{8}{5}$ and $FE = FC + CE = \frac{16}{5} + \frac{4}{3} = \frac{68}{15}$, so by the Pythagorean Theorem, $AE = \dfrac{20\sqrt{13}}{5}$. You can now use similar triangles to find that $x =  \dfrac{8}{5\sqrt{13}}$ and continue on like in solution 2.

 





\textbf{Solution 4}

$DE$ goes through $A$, the point of tangency of both circles. So $DE$ intercepts equal arcs in circle $P$ and $Q$:  \href{https://artofproblemsolving.com/wiki/index.php?title=Homothety}{homothety}. Hence, $AE=4AD$. We will use such similarity later.

 The diagonal distance between the centers of the circles is $4+1=5$. The difference in heights is $4-1=3$. So $BC=\sqrt{5^2-3^2}=4$.

 The triangle connecting the centers with a side parallel to $BC$  is a $3-4-5$ right triangle. Since $O_PA=1$, the height of $A$ is $1+3/5=8/5$. Drop an altitude from $A$ to $BC$ and call it $I$: $IB=4/5$ and $IC=4-4/5=16/5$. Since right $\triangle AIB\sim\triangle CIB$, $ABC$ is a right triangle also; $IB:IA:IC$ form a geometric progression $\times 2$.

 
Extend $BA$ through $A$ to a point $G$ on the other side of $\circ Q$. By \href{https://artofproblemsolving.com/wiki/index.php?title=Homothety}{homothety}, $\triangle DAB\sim\triangle EAG$. By angle chasing $\triangle DAB$ through right triangle $ABC$, we deduce that $\angle CEG$ is a right angle. Since $ACEG$ is cyclic, $\angle GAC$ is also right. So $CG$ is a diameter of $\circ G$. Because of this, $CG \perp BC$, the tangent line. $\triangle BCG$ is right and $\triangle BCG\sim\triangle ABC\sim\triangle CAG$.

 
$AC=\sqrt{(8/5)^2+(16/5)^2}=8\sqrt{5}/5$ so $AG=2AC=16\sqrt{5}/5$ and $[\triangle CAG]=64/5$. 

 Since $[\triangle DAB]=[\triangle ACE]$, the common area is $[ACEG]/17$. $16[\triangle DAB]=[\triangle GAE]$ because the triangles are similar with a ratio of $1:4$. So we only need to find $[\triangle CEG]$ now.

 
Extend $DE$ through $E$ to intersect the tangent at $F$. Because $4DA=AE$, the altitude from $B$ to $AD$ is $4$ times the height from $C$ to $EA$. So $BC=3/4BF$ and $BF=16/3$. We look at right triangle $\triangle AIF$. $IF=68/15$ and $AI=8/5$. $\triangle AIF$ is a $17-6-5\sqrt{13}$ right triangle. Hypotenuse $AF$ intersects $CG$ at a point, we call it $H$. $CH=4/3\div 68/15\cdot 8/5=8/17$. So $HG=8-8/17=128/17$.

 
By \href{https://artofproblemsolving.com/wiki/index.php?title=Power\_of\_a\_Point\_Theorem}{Power of a Point}, $CH\cdot HG=AH\cdot HE$. $AH=16/5\cdot 5\sqrt{13}/17=16\sqrt{13}/17.$ So $HE=1024/289\cdot 17/(16\sqrt{13})=64/(17\sqrt{13})$. The height from $E$ to $CG$ is $17/(5\sqrt{13})\cdot 64/(17\sqrt{13})=64/65$.

 
Thus, $[\triangle CEG]=64/65\cdot 8\div 2=256/65$. The area of the whole cyclic quadrilateral is $64/5+256/65=(832+256)/65=1088/65$. Lastly, the common area is $1/17$ the area of the quadrilateral, or $64/65$. So $64+65=\boxed{129}$.

 





\textbf{Solution 5 (HARD computation)}


                    \begin{figure}[H]
                        \centering
                        \includegraphics[width=0.5\textwidth]{Solution/AIME_2015_II_15_Solution_Figure_1.png}
                    \end{figure}
                    

 Consider the homothety that takes triangle BDA onto CXY on the big circle, as plotted. Some hidden congruence angles are revealed which help reduce computation complexity. Just some angle chasing and straight forward trigs. Because $AE=XY$ and $AE \parallel XY$, $XYE$ is right angle.  

 First, $\frac{NO_1}{NO_1+5} = \frac{1}{4}$, so $NO_1=\frac{5}{3}$. And,
 \[\cos{\angle{AO_2C}}=\cos{2\angle{AYC}} = \frac{O_2C}{NO_1+5} = \frac{3}{5}\]
 \[\sin{\angle{AYC}} = \sqrt{\dfrac{1-\cos{2\angle{AYC}}}{2}}=\frac{1}{\sqrt{5}}\]
 \[\cos{\angle{AYC}} = \frac{2}{\sqrt{5}}\]Then,
 \[[AEC] = \frac{1}{2}AE*CE*\sin{\angle{ACE}}=\frac{1}{2}AE*8\sin{\angle{CYE}}*\frac{1}{\sqrt{5}}\]
 \[[CXY] = \frac{1}{2}CX*XY*\sin{\angle{CXY}}=\frac{1}{2}*8\sin{\angle{XYC}}*XY*\sin{\angle{CAY}}\]Since $\angle{CAY} = 90 - \angle{AYC}$, $\angle{XYC} = 90 - \angle{CYE}$, $XY = AE$, we have
 \[[CXY] = \frac{1}{2}*8AE\cos{\angle{CYE}}*\cos{\angle{AYC}}=\frac{1}{2}*8AE\cos{\angle{CYE}}*\frac{2}{\sqrt{5}}\]Since $\triangle{CXY}$ is four times in scale to $\triangle{AEC}$, their area ratio is 16. Divide the two equations for the two areas, we have
 \[\tan{\angle{CYE}} = \frac{1}{8}\]With this angle found, everything else just follows. 
 \[\sin{\angle{CYE}} = \dfrac{\tan{\angle{CYE}}}{\sqrt{1+\tan^2{\angle{CYE}}}}=\dfrac{1}{\sqrt{65}}\]
 \[\cos{\angle{CYE}} = \dfrac{8}{\sqrt{65}}\]
 \[\sin{\angle{AYE}} = \sin({\angle{AYC}+\angle{CYE}} )= \frac{1}{\sqrt{5}}*\frac{8}{\sqrt{65}} + \frac{2}{\sqrt{5}}*\frac{1}{\sqrt{65}} = \frac{2}{\sqrt{13}}\]
 \[AE = 8\sin{\angle{AYE}} = \frac{16}{\sqrt{13}}\]
 \[[AEC] = \frac{1}{2}*8*\frac{16}{\sqrt{13}}*\dfrac{1}{\sqrt{65}}*\frac{1}{\sqrt{5}}=\frac{64}{65}\]Thus, our answer is $64+65=\boxed{129}$.

 





\textbf{Solution 6 (Simple computation)}

Let $K$ be the intersection of $BC$ and $AE$. Since the radii of the two circles are 1:4, so we have $AD:AE=1:4$, and the distance from $B$ to line $l$ and the distance from $C$ to line $l$ are in a ratio of 4:1, so $BK:CK=4:1$. We can easily calculate the length of $BC$ to be 4, so $CK=\frac{4}{3}$. Let $J$ be the foot of perpendicular line from $A$ to $BC$, we can know that $BJ:CJ=1:4$, so $BJ = 0.8$, $CJ=3.2$, $AJ=1.6$, and $AK=\sqrt{1.6^2+\left(3.2+\frac{4}{3}\right)^2}=\frac{4}{3}\sqrt{13}$. Since $CK^2 = EK\cdot AK$, so $EK=\frac{4}{39}\sqrt{13}$, and $AE = \frac{4}{3}\sqrt{13} - \frac{4}{39}\sqrt{13}= \frac{16}{13}\sqrt{13}$. $\sin\angle AKB=\frac{AJ}{AK} = \frac{1.6}{\frac{4}{3}\sqrt{13}}=\frac{1.2}{\sqrt{13}}$, so the distance from $C$ to line $l$ is $d=CK\cdot \sin\angle AKB = \frac{4}{3}\cdot \frac{1.2}{\sqrt{13}}=\frac{1.6}{\sqrt{13}}$. so the area is
 \[[ACE] = \frac{1}{2}\cdot AE\cdot d = \frac{1}{2}\cdot\frac{16}{13}\sqrt{13}\frac{1.6}{\sqrt{13}} = \frac{64}{65}\]The final answer is $\boxed{129}$.

 --- by Dan Li

 





\textbf{Solution 7}

Consider the common tangent from $A$ to both circles. Let this intersect $BC$ at point $K$. From equal tangents, we have $BK=AK=CK$, which implies that $\angle BAC = 90^\circ$.

 Let the center of $\mathcal{P}$ be $O_1$, and the center of $\mathcal{Q}$ be $O_2$. Angle chasing, we find that $\triangle O_1DA \sim \triangle O_2EA$ with a ratio of $1:4$. Hence $4AD = AE$.

 We can easily deduce that $BC=4$ by dropping an altitude from $O_1$ to $O_2C$. Let $\angle ABC = \theta$. By some simple angle chasing, we obtain that $\angle BO_1A = 2\angle BDA = 2\angle ABC = 2\theta,$ and similarly $\angle CO_2A = 180 - 2\theta$.

 Using LoC, we get that $AB = \sqrt{2-2\cos2\theta}$ and $AC = \sqrt{32+32\cos2\theta}$. From Pythagorean theorem, we have 
 \[AB^2 + AC^2 = BC^2 \implies \cos 2\theta = -\frac{3}{5} \implies \cos \theta = \frac{1}{\sqrt5}, \sin \theta = \frac{2}{\sqrt 5}\]In other words, $AB = \frac{4}{\sqrt5}, AC = \frac{8}{\sqrt5}$.

 Using the area condition, we have:
 \begin{align*} \frac12 AD*AB \sin \angle DAB &= \frac 12 AE*AC \sin(90-\angle DAB) \\ AD*\frac{4}{\sqrt5} \sin \angle DAB &= 4AD*\frac{8}{\sqrt5} \cos \angle DAB \\ \sin \angle DAB &= 8 \cos \angle DAB \\ \implies \sin \angle DAB &= \frac{8}{\sqrt{65}} \end{align*}


 Now, for brevity, let $\angle D = \angle ADB$ and $\angle A = \angle DAB$.

 From Law of Sines on $\triangle ABD$, we have
 \begin{align*} \frac{AB}{\sin \angle D} &= \frac{AD}{\sin (180-\angle A - \angle D)} \\ \frac{\frac{4}{\sqrt5}}{\frac{2}{\sqrt5}} &= \frac{AD}{\sin \angle A \cos\angle D + \sin \angle D\cos\angle A} \\ 2 &= \frac{AD}{\frac{2}{\sqrt{13}}} \\ AD &= \frac{4}{\sqrt{13}} \end{align*}


 It remains to find the area of $\triangle ABD$. This is just 
 \[\frac12 AD*AB*\sin \angle A = \frac12 * \frac{4}{\sqrt{13}}*\frac{4}{\sqrt5}*\frac{8}{\sqrt{65}} = \frac{64}{65}\] for an answer of $\boxed{129}.$

 This solution was brought to you by Leonard\_my\_dude.

 





\textbf{Solution 8 (Synthetic-Trigonometry)}

Add in the line $k$ as the internal tangent between the two circles. Let $M$ be the midpoint of $BC$; It is well-known that $M$ is on $k$ and because $k$ is the radical axis of the two circles, $AM=BM=CM$. Therefore because $M$ is the circumcenter of $\triangle{BAC}$, $\angle{BAC}=90^{\circ}$. Let $O_P$ be the center of circle $\mathcal{P}$ and likewise let $O_Q$ be the center of circle $\mathcal{Q}$. It is well known that by homothety $O_P, A,$ and $O_Q$ are collinear. It is well-known that $\angle{ABC}=\angle{ADB}=b$, and likewise $\angle{ACB}=\angle{AEC}=c$. By homothety, $AD=4AE$, therefore since the two triangles mentioned in the problem, the length of the altitude from $B$ to $AD$ is four times the length of the altitude from $C$ to $AE$. Using the Pythagorean Theorem, $BC = 4$. By angle-chasing, $O_{P}AMB$ is cyclic, and likewise $O_{Q}CMA$ is cyclic. Use the Pythagorean Theorem for $\triangle{O_{P}BM}$ to get $O_{P}M=\sqrt{5}$. Then by Ptolemy's Theorem $AB=\frac{4\sqrt{5}}{5} \implies AC=\frac{8\sqrt{5}}{5}$. Now to compute the area, using what we know about the length of the altitude from $B$ to $AD$ is four times the length of the altitude from $C$ to $AE$, letting $x$ be the length of the altitude from $C$ to $AE$, $x=\frac{8}{5\sqrt{13}}$. From the Law of Sines, $\frac{BD}{\sin{A}}=2 \implies \sin{A}=\frac{4x}{AB} \implies BD=\frac{8x}{AB}=\frac{16\sqrt{65}}{65}$. Then use the Pythagorean Theorem twice and add up the lengths to get $AD=\frac{4}{\sqrt{13}}$. Use the formula $\frac{b \times h}{2}$ to get $\frac{64}{65} = \boxed{129}$ as the answer.

 





\textbf{Solution 9 (Visual)}


                    \begin{figure}[H]
                        \centering
                        \includegraphics[width=0.5\textwidth]{Solution/AIME_2015_II_15_Solution_Figure_1.png}
                    \end{figure}
                    
                    \begin{figure}[H]
                        \centering
                        \includegraphics[width=0.5\textwidth]{Solution/AIME_2015_II_15_Solution_Figure_1.png}
                    \end{figure}
                    Let $P$ and $Q$ be the centers of circles $\mathcal{P}$ and $\mathcal{Q}$ , respectively.

 Let $M$ be midpoint $BC, \beta = \angle ACB.$ 

 Upper diagram shows that

 $\sin 2\beta = \frac {4}{5}$ and $AC = 2 AB.$ Therefore $\cos 2\beta = \frac {3}{5}.$

 Let $CH\perp l, BH'\perp l.$ Lower diagram shows that 

 $\angle CAE = \angle ABH' = \alpha$ (perpendicular sides) 

 and $\angle CQE = 2\alpha$ (the same intersept $\overset{\Large\frown} {CE}).$
 \[\tan\alpha = \frac {1}{8}, \sin2\alpha = \frac{2 \tan \alpha}{1 + \tan^2 \alpha} = \frac {16}{65}, \cos2\alpha = \frac{1 - \tan^2 \alpha}{1 + \tan^2 \alpha} = \frac {63}{65}.\]The area
 \[[ACE] = [AQC]+[CQE]– [AQE].\]Hence 
 \[[ACE] =\frac{AQ^2}{2} \left(\sin 2\alpha + \sin 2\beta - \sin(2\alpha + 2\beta)\right),\]
 \[[ACE]  = 8\left( \frac{16}{65}+\frac{4}{5} - \frac{4}{5}\cdot \frac{63}{65} - \frac{3}{5}\cdot \frac{16}{65}\right) = \frac{64}{65}\implies \boxed {129}.\]\textbf{vladimir.shelomovskii@gmail.com, vvsss}

 





\textbf{Solution 10 (Similar Triangles, Angle Chasing, and Ptolemy's)}

We begin by extending $\overline{AB}$ upwards until it intersects Circle $\mathcal{Q}$. We can call this point of intersection $F$. Connect $F$ with $E$, $C$, and $A$ for future use.

 Create a trapezoid with points $B$, $C$, and the origins of Circles $\mathcal{P}$ and $\mathcal{Q}$. After quick inspection, we can conclude that the distance between the origins is 5 and that $\overline{BC}$ is 4.

 (Note: It is important to understand that we could simply use the formula for the distance between the tangent points on a line and two different circles, which is $2\cdot \sqrt{a\cdot b}$, where $a$ and $b$ are the two respective radii of the circles. In our case, we get $2\cdot \sqrt{4} = 4$.

 Using similar triangles or homotheties, $AE=4\cdot AD$ and $AF=4\cdot BA$. $BC^2 = BA\cdot BF$. 
 \[16 = BA\cdot (5\cdot BA)\] 
 \[BA = \frac{4}{\sqrt{5}}\] 
 \[AF = \frac{16}{\sqrt{5}}\] 
 \[BF = 4\cdot \sqrt{5}\]

 Inspecting $\triangle{BFC}$, we recognize that it is a right triangle ($\angle{BCF} = 90$) as the final length ($\overline{FC}$) being 8 would allow for an $x-2x-x\sqrt{5}$ triangle. Hence, the diameter of circle $\mathcal{Q}$ = $CF$. This also means that $\angle{CAF} = \angle{CEF} = 90$.

 From the fact that $\triangle{ABC}$ is a right triangle: 
 \[AB^2 + AC^2 = BC^2\] 
 \[AC = \frac{8}{\sqrt{5}}\](Note: We could have also  used $\triangle{FAC}$.)

 Our next step is to start angle chasing to find any other similar triangles or shared angles. Label the intersection between $\overline{FC}$ and $\overline{DE}$ as $M$. Label $\angle{CAE} = \angle{CFE} = \theta$. Since $\angle{BAC} = 90$, $\angle{EAF} = \angle{DAB} = 90-\theta$. Now, use the sine formula and the fact that the areas of $\triangle{DBA}$ and $\triangle{ACE}$ are equal to get:
 \[\frac{1}{2}\cdot AC\cdot AE\cdot \sin{\theta} = \frac{1}{2}\cdot AD\cdot AB\cdot \sin{(90-\theta)}\]Since $\sin{(90-\theta)} = \cos{\theta}$:
 \[\tan{\theta} = \frac{AD}{2\cdot AE}\]

 Using right $\triangle{FCE}$, since $\angle{ACM} = \angle{CFE}$, $\tan{\theta} = \frac{CE}{FE}$. Hence, plugging into the previous equation:
 \[\frac{CE}{FE} = \frac{AD}{2\cdot AE}\]Using the Pythagorean theorem on $\triangle{FCE}$, $FE = \sqrt{(64 - CE^2)}$. We also know that $AE=4\cdot AD$Plugging back in:
 \[\frac{CE}{\sqrt{(64 - CE^2)}} = \frac{AD}{2\cdot (4\cdot AD)}\]
 \[\frac{CE}{\sqrt{(64 - CE^2)}} = \frac{1}{8}\].From here, we can square both sides and bring everything to one side to get:
 \[EF^2 + 64EF - 64 = 0\].
 \[EF = \frac{64}{\sqrt{65}}\]
 \[CE = \frac{8}{\sqrt{65}}\]We should also return to the fact that $\sin{\theta} = \frac{CE}{CF}$ from $\triangle{FCE}$, so 
 \[\sin{\theta} = \frac{1}{\sqrt{65}}\]

 From the fact that $\angle{CAF} = \angle{CEF} = 90$, we can use Ptolemy's Theorem on quadrilateral $ACEF$. $AC\cdot EF + CE\cdot FA = CF\cdot AE$. Plugging in and solving, we get that $AE = \frac{16}{\sqrt{13}}$.

 We now have all of our pieces to use the Sine Formula on $\triangle{ACE}$. 
 \[\frac{1}{2}\cdot AC\cdot AF\cdot \sin{\theta}\]
 \[\frac{1}{2}\cdot \frac{8}{\sqrt{5}}\cdot \frac{16}{\sqrt{13}}\cdot \frac{1}{\sqrt{65}} = \frac{64}{65} = \boxed {129}\]

 





\textbf{Solution 11 (Trig)}

Let the center of the larger circle be $O,$ and the center of the smaller circle be $P.$ It is not hard to find the areas of $ACO$ and $ABP$ using pythagorean theorem, which are $\frac{32}{5}$ and $\frac{2}{5}$ respectively. Assign $\angle AOC=a,\angle COE=c,\angle DPB=d,\angle BPA=b.$ We can figure out that $\angle AOE=\angle DOA=\theta$ using vertical angles and isosceles triangles. Now, using $[ABC]=\frac{1}{2}ab\sin C$
 \[[ACE]=[AOC]+[COE]-[AOC]=\dfrac{32}{5}-8\sin c-8\sin \theta,\]
 \[[DAB]=[APD]+[APB]+[BPD]=\dfrac{2}{5}+\dfrac{1}{2}\sin \theta+\dfrac{1}{2}\sin d.\]We can also figure out that $\sin a=\dfrac{4}{5},\cos a=\dfrac{3}{5},\sin b=\dfrac{4}{5}, \cos b=-\dfrac{3}{5}.$ Also, $c=\theta-a$ and $d=360-\theta-b.$ Using sum and difference identities:
 \[\sin c=\dfrac{3}{5}\sin \theta-\dfrac{4}{5}\cos \theta,\]
 \[\sin d=\dfrac{3}{5}\sin\theta-\dfrac{4}{5}\cos\theta.\](We can also notice that $c+d=360-\theta-b+\theta-a=360-(a+b)=180$ which means that $\sin c=\sin d.$)Substituting in the equations for $\sin c$ and $\sin d$ into the equations for $[ACE]$ and $[DAB],$ setting them equal, and simplifying: 
 \[3=2\sin\theta+3\cos\theta.\]Solving this equation we get that $\sin\theta=\frac{12}{13}$ and $\cos\theta=\frac{5}{13}.$ Doing a lot of substitution gives us 
 \[[ACE]=[DAB]=\dfrac{64}{65},\] which means the answer is $64+65=\boxed{129}.$

 





\textbf{Solution 12 (Coordinates and Shoelace)}


                    \begin{figure}[H]
                        \centering
                        \includegraphics[width=0.5\textwidth]{Solution/AIME_2015_II_15_Solution_Figure_1.png}
                    \end{figure}
                    

 Let the centers of circles $\mathcal{P}$ and $\mathcal{Q}$ be $P$ and $Q$, respectively. Let $F$ be the point on $\overline{QC}$ such that $\overline{PF}$ is parallel to $\overline{BC}$. We find that $BC=PF=\sqrt{(4+1)^{2}-(4-1)^2}=4$.

 Now, we define $A$ to be $(0,0)$. It follows that:

 \begin{align}P &= \left(-\frac{4}{5},-\frac{3}{5}\right) \\Q &= \left(\frac{16}{5},\frac{12}{5}\right) \\B &= \left(-\frac{4}{5},-\frac{8}{5}\right) \\C &= \left(\frac{16}{5},-\frac{8}{5}\right) \\\end{align}

 The equation of circle $P$ is $\left(x+\frac{4}{5}\right)^2+\left(y+\frac{3}{5}\right)^2 = 1^2$, which simplifies to $x^2+\frac{8}{5}x+y^2+\frac{6}{5}y=0$.

 The equation of circle $Q$ is $\left(x-\frac{16}{5}\right)^2+\left(y-\frac{12}{5}\right)^2 = 4^2$, which simplifies to $x^2-\frac{32}{5}x+y^2-\frac{24}{5}y=0$.

 Suppose points $D$ and $E$ lie on the line $y=mx$ for some slope $m$. We can now plug in $y=mx$ into the equations for the circles to find the coordinates of $D$ and $E$ in terms of $m$.

 For point $D$:
 \[x^2+\frac{8}{5}x+m^2x^2+\frac{6}{5}mx=0\]
 \[(1+m^2)x^2+\left(\frac{8}{5}+\frac{6}{5}m\right)=0\]Since $D \neq A$, we have $x \neq 0$, thus $x=\dfrac{-\left(\frac{8}{5}+\frac{6}{5}m\right)}{1+m^2}$.

 For point $E$:
 \[x^2-\frac{32}{5}x+m^2x^2-\frac{24}{5}mx=0\]
 \[(1+m^2)x^2-\left(\frac{32}{5}+\frac{24}{5}m\right)=0\]Since $E \neq A$, we have $x \neq 0$, thus $x=\dfrac{\left(\frac{32}{5}+\frac{24}{5}m\right)}{1+m^2}$.

 We let $a=\dfrac{\left(\frac{8}{5}+\frac{6}{5}m\right)}{1+m^2}$ to write things quicker, and so the coordinates of $D$ are $(-a,-am)$ and the coordinates of $E$ are $(4a,4am)$. 

 Now, we can apply shoelace on $\triangle DBA$ and $\triangle ACE$:$[DBA] = \frac {1}{2}\left|(0 \cdot -\frac{8}{5} + -\frac{4}{5} \cdot -am + -a \cdot 0) - (0 \cdot -\frac{4}{5} + -\frac{8}{5} \cdot -a + -am \cdot 0)\right| = \dfrac{\left|-\frac{4}{5}am + \frac{8}{5}a\right|}{2} = \dfrac{\frac{8}{5}a-\frac{4}{5}am}{2}$.

 $[ACE] = \frac {1}{2}\left|(0 \cdot -\frac{8}{5} + \frac{16}{5} \cdot 4am + 4a \cdot 0) - (0 \cdot \frac{16}{5} + -\frac{8}{5} \cdot 4a + 4am \cdot 0)\right| = \dfrac{\left|-\frac{32}{5}a-\frac{64}{5}am\right|}{2} = \dfrac{\frac{32}{5}a+\frac{64}{5}am}{2}$. 

 Note on determining the absolute value: We can assume $-\frac{1}{2}<m<0$ because $\overline{CA}$ has a slope of $-\frac{1}{2}$ and the area of $\triangle ACE$ gets too big if $m>0$. So, we can also assume $a$ is positive.

 Equating the two areas yields $\dfrac{\frac{8}{5}a-\frac{4}{5}am}{2} = \dfrac{\frac{32}{5}a+\frac{64}{5}am}{2} \implies -\frac{24}{5}a=\frac{68}{5}am \implies m=-\frac{6}{17}$.

 We also know that $a=\dfrac{\left(\frac{8}{5}+\frac{6}{5}m\right)}{1+m^2}=\dfrac{\frac{8}{5}-\frac{36}{85}}{1+\frac{36}{289}}=\dfrac{\frac{100}{85}}{\frac{325}{289}}=\frac{68}{65}$.

 Now, we can compute the area of $\triangle DBA$: $[DBA] = \dfrac{\frac{8}{5} \cdot \frac{68}{65} -\frac{4}{5} \cdot \frac{68}{65} \cdot \left(-\frac{6}{17}\right)}{2}=\dfrac{\frac{8\cdot68+4\cdot4\cdot6}{5\cdot65}}{2}=\frac{64}{65} \implies m+m=\boxed{129}$.

 





\textbf{Solution 13 (Only Similar Triangles, No Trig)}


                    \begin{figure}[H]
                        \centering
                        \includegraphics[width=0.5\textwidth]{Solution/AIME_2015_II_15_Solution_Figure_1.png}
                    \end{figure}
                    Let $BH=x$, $HC = 4-x$

 $x(4-x) = (\frac{8}{5})^2$

 $x^2-4x+\frac{64}{25} = 0$

 $25x^2-100x+64 - 0$

 $(5x-16)(5x-4) = 0$

 So $HC = \frac{16}{5}$, $BH = \frac{4}{5}$

 
$BC = \sqrt{(1+4)^2-(4-1)^2} = 4$

 $\frac{r_1}{r_2} = \frac{AD}{AE} = \frac{1}{4} => \frac{h_1}{h_2} = 4$

 Because $\triangle QIC \sim \triangle QGB$

 $\frac{QC}{QB} = \frac{1}{4}$

 $4QC = QB = QC+4$

 $QC = \frac{4}{3}$

 Because $O_1A = \frac{1}{5} O_2$

 $AH = \frac{4}{5} O_1B = \frac{1}{5} O_2C = \frac{4}{5}+\frac{4}{5} = \frac{8}{5}$

 $AQ = \sqrt{(\frac{8}{5})^2+(\frac{68}{15})^2} = \frac{4}{3}\sqrt{13}$

 $CQ^2 = EQ\cdot AQ => EQ = \frac{\frac{16}{9}}{\frac{4}{3}\sqrt{13}} = \frac{4}{39}\sqrt{13}$

 $AE = \frac{4}{3}\sqrt{13}-\frac{4}{39}\sqrt{13} = \frac{16}{13}\sqrt{13}$

 $\sin{Q} = \frac{AH}{AQ} = \frac{\frac{8}{5}}{\frac{4}{3}\sqrt{13}} = \frac{6}{5\sqrt{13}}$

 $CI = CQ\cdot \sin{Q} = \frac{4}{3}\frac{6}{5\sqrt{13}} = \frac{8}{5\sqrt{13}}$

 So $S_{shaded} = \frac{1}{2}\cdot \frac{16}{13}\sqrt{13}\cdot \frac{8}{5\sqrt{13}} = \frac{64}{65}$

 $64+65 = \boxed{129}$

 -cassphe

 





\textbf{Solution 14 (Coordinate Geometry)}

Let $O_1, O_2$ be the center of circles $\mathcal{P},\mathcal{Q}$, respectively. Note that $A=O_1+\frac{1}{5}\overrightarrow{O_{1}O_{2}}$. Let $X$ be the intersection of $\overline{DE}$ and $\overline{BC}$. We establish a coordinate system with $B$ being the origin and $\overline{BC}$ being the $x$-axis. After some minor calculations, here is our setup:\begin{align*}A&=(\frac{4}{5},\frac{8}{5})\\B&=(0,0)\\C&=(4,0)\\O_1&=(0,1)\\O_2&=(4,4)\\X&=(x_1,0)\\D&=(x_2,y_2)\\E&=(x_3,y_3)\\DE&: y=-k(x-\frac{4}{5})+\frac{8}{5}\\\mathcal{P}&: x^2+(y-1)^2=1\\\mathcal{Q}&: (x-4)^2+(y-4)^2=16\end{align*}
We see that the equation of $DE$ can also be written as $y=-k(x-x_1)$. Equating the constant term, we find
 \[x_1=\frac{4k+8}{5k}.\]The equal-area condition implies\begin{equation}\frac{h_1}{h_2}\cdot \frac{AE}{AD}=1 \tag{1}\end{equation}where $h_1$ is the altitude from $C$ to $\overline{DE}$ and $h_2$ is the altitude from $B$ to $\overline{DE}$.By similarity, we compute $\frac{h_1}{h_2}$:\begin{align*}\frac{h_1}{h_2}&=\frac{CX}{BX}\\\frac{h_1}{h_2}&=\frac{x_{1}-4}{x_{1}}\\\frac{h_1}{h_2}&=\frac{2-4k}{k+2}\end{align*}
Solving the system of equation of $DE$ and $\mathcal{P}$, we have
 \[(1+k^2)x^2-(\frac{8}{5}k^2+\frac{6}{5}k)x+C=0.\]By the Vieta's formula,\begin{equation}\frac{4}{5}+x_2=\frac{\frac{8}{5}k^2+\frac{6}{5}k}{1+k^2}\implies x_2=\frac{4k^2+6k-4}{5+5k^2} \tag{2}\end{equation}Computing the $x$-coordinate of $E$ similarly, we find
 \[x_3=\frac{4k^2-24k+36}{5+5k^2}.\]We proceed by computing $\frac{AE}{AD}$:\begin{align*}\frac{AE}{AD}&=\frac{x_3-\frac{4}{5}}{\frac{4}{5}-x_2}\\\frac{AE}{AD}&=\frac{\frac{-24k+32}{5+5k^2}}{\frac{-6k+8}{5+5k^2}}\\\frac{AE}{AD}&=4\end{align*}
Solving $(1)$, we have
 \[4\cdot\frac{2-4k}{k+2}=1\implies k=\frac{6}{17}\]By $(2)$ and the equation of $DE$, we find
 \[D=(-\frac{16}{65},\frac{128}{65})\]Using the Shoelace Formula, we obtain the area $[BAD]$:\begin{align*}[BAD]&=\frac{1}{2}\lvert -\frac{16}{65}\cdot\frac{8}{5}+\frac{4}{5}(0-\frac{128}{65}) \rvert\\[BAD]&=\frac{1}{2}\cdot\frac{16\cdot 8+4\cdot 128}{65\cdot 5}\\[BAD]&=\frac{64}{65}\end{align*}
The requested sum is $64+65=\boxed{129}$.

 




\end{enumerate}